
\documentclass[aps,prx,reprint,superscriptaddress,longbibliography]{revtex4-2}

\usepackage{amsmath,amssymb,amsthm}
\usepackage{bm}
\usepackage{microtype}
\usepackage{graphicx}
\usepackage{hyperref}
\usepackage{booktabs}
\usepackage{multirow}
\usepackage{xcolor}
\usepackage{cleveref}
\usepackage{enumitem}
\usepackage{thmtools}
\usepackage{pifont}

\hypersetup{
    colorlinks=true,
    linkcolor=blue,
    citecolor=blue,
    urlcolor=blue
}

\emergencystretch=1.5em

\newcommand{\ket}[1]{|#1\rangle}
\newcommand{\bra}[1]{\langle#1|}
\newcommand{\braket}[2]{\langle#1|#2\rangle}

\newcommand{\QAOA}{\textnormal{\textsc{QAOA}}}
\newcommand{\GNN}{\textnormal{\textsc{GNN}}}
\newcommand{\GIN}{\textnormal{\textsc{GIN}}}
\newcommand{\UQQAOA}{\textnormal{\textsc{UQ-QAOA}}}
\newcommand{\Normal}{\mathcal{N}}
\newcommand{\calG}{\mathcal{G}}
\newcommand{\calV}{\mathcal{V}}
\newcommand{\calE}{\mathcal{E}}
\newcommand{\calD}{\mathcal{D}}
\newcommand{\calN}{\mathcal{N}}
\newcommand{\calL}{\mathcal{L}}
\newcommand{\calF}{\mathcal{F}}
\newcommand{\calH}{\mathcal{H}}
\newcommand{\calP}{\mathcal{P}}
\newcommand{\calB}{\mathcal{B}}
\newcommand{\calT}{\mathcal{T}}
\newcommand{\calS}{\mathcal{S}}
\newcommand{\calR}{\mathcal{R}}
\newcommand{\bbR}{\mathbb{R}}
\newcommand{\bbE}{\mathbb{E}}
\newcommand{\bbP}{\mathbb{P}}
\newcommand{\bbN}{\mathbb{N}}

\newtheorem{theorem}{Theorem}
\newtheorem{proposition}[theorem]{Proposition}
\newtheorem{corollary}[theorem]{Corollary}
\newtheorem{lemma}[theorem]{Lemma}
\newtheorem{remark}{Remark}
\newtheorem{definition}{Definition}

% ============================================================
\begin{document}
% ============================================================

\title{Graph-Conditioned Mahalanobis Trust Regions for Query-Efficient \QAOA{}}

\author{Molena Huynh}
\email{Contact author: Molena.Huynh@jmp.com}
\affiliation{North Carolina State University}

\date{}

% ============================================================
\begin{abstract}
In low-depth variational quantum optimization, the dominant resource
is often the number of circuit evaluations rather than the classical
postprocessing surrounding them: every candidate \QAOA{} angle vector
must be estimated from measurements. Existing transfer methods predict
a single warm start and then discard uncertainty after the first
iterate. We instead learn a graph-conditioned diagonal Gaussian
$\Normal(\bm\mu,\bm\Sigma)$ over angles and reuse it throughout
inference: the mean initializes local search, the covariance defines a
Mahalanobis trust region that constrains every iterate, and the trace
of the covariance sets an instance-adaptive evaluation budget. The
predictor is a Graph Isomorphism Network with Laplacian spectral
encodings, trained with negative log-likelihood, a Wasserstein
regularizer that couples nearby instances in angle space, and a
supervised contrastive loss that orders embeddings by landscape
similarity before the Gaussian heads act. We prove local guarantees
for trust-region confinement, conditional gradient
anti-concentration, preservation of expected ordering under affine
depolarizing noise, spectral-norm generalization, and finite-sample
conformal coverage. On depth-$2$ MaxCut across Erd\H{o}s--R\'enyi,
3-regular, Barab\'asi--Albert, and Watts--Strogatz graphs with
$n{=}8$--$16$, training only at $n{=}14$, the mean evaluation count at
$n{=}14$ drops to $45\pm 7$ from $343$ for random restarts and $85$
for the strongest learned-point baseline, an $86.9\%$ reduction,
while the sampled best-bitstring approximation ratio remains within
$0.027$ of the best concentration-based heuristic; at a higher-budget
operating point the gap shrinks to $0.007$. Uncertainty is calibrated
($\mathrm{ECE}{=}0.052$, Spearman $\rho{=}0.770$), the speedup remains
$7.2$--$7.7\times$ under finite-shot noise, and the gains transfer to
unseen graph sizes and graph families without retraining.
\end{abstract}

\maketitle

\section{Introduction}
\label{sec:intro}

The Quantum Approximate Optimization Algorithm
(\QAOA)~\cite{farhi2014} applies $p$ alternating layers of a cost
unitary and a mixer unitary to the uniform superposition.  The $2p$
angles are chosen to maximize the expected value of the cost
operator.  At low depth, derivative-free solvers often require
hundreds of objective evaluations per instance~\cite{zhou2020,wurtz2021}.
On quantum hardware, each evaluation requires repeated state
preparation and measurement.  Query count is therefore a primary
resource~\cite{cerezo2021,preskill2018}.

We reduce this cost by predicting a distribution rather than a point.
A Graph Isomorphism Network (\GIN)~\cite{xu2019} with Laplacian
spectral encodings~\cite{dwivedi2023,lim2023} maps a graph $\calG$ to
a Gaussian $\Normal(\bm\mu(\calG),\bm\Sigma(\calG))$ over \QAOA{}
angles.  The mean initializes local optimization.  The covariance
defines the region to which the optimizer is confined.  The trace of
the covariance sets the instance-dependent evaluation budget.

This differs from a warm start.  A warm start uses the predictor once;
here the covariance constrains every iterate and also sets the budget.

The scope is restricted to query efficiency in a small-graph,
low-depth regime where exact statevector simulation and finite-shot
simulation are feasible on a CPU.

The empirical study isolates four questions: whether the method
reduces query count at competitive quality, whether the gain comes
from uncertainty-aware search rather than initialization alone,
whether it transfers across graph sizes and families, and whether the
predicted uncertainty is calibrated enough to drive allocation.

Section~\ref{sec:theory} gives the local guarantees: control of the
loss induced by confinement, a conditional anti-concentration
statement for the gradient, preservation of expected ordering under
depolarizing noise, and finite-sample coverage after conformal
calibration.

We evaluate the method on MaxCut at depth $p{=}2$.  The test graphs
are Erd\H{o}s--R\'enyi, 3-regular, Barab\'asi--Albert, and
Watts--Strogatz families with $n{=}8$--$16$ vertices.  Training uses
240 graphs at $n{=}14$ only.  The learned trust region transfers to
families and sizes not seen during training.  On $n{=}14$ graphs the
mean evaluation count drops from $343$ for random restarts and $85$
for the best learned-point predictor to $45\pm 7$, while sampled
approximation ratios remain within $3$ percentage points of the
concentration heuristic.  The factor-of-two gain over the learned
point predictor is the effect of the trust-region constraint and the
adaptive budget.  It is not the effect of initialization alone.

The main contributions of this work are:
\begin{enumerate}[leftmargin=*,itemsep=2pt]
\item \textbf{Gaussian angle predictor.}
A \GIN{} with Laplacian spectral positional encodings trained to
predict a full diagonal Gaussian $\Normal(\bm\mu,\bm\Sigma)$ over
\QAOA{} angles, where both the mean head and the variance head are
jointly trained through a two-phase objective combining NLL,
Wasserstein regularization, and supervised contrastive loss.
\item \textbf{Mahalanobis trust region.}
A Mahalanobis ellipsoid $\calT_\alpha(\calG)$ constructed from the
predicted covariance and used as a hard constraint throughout the
Nelder--Mead refinement, including projection of iterate that exit the
region.  The geometry is $\bm\Sigma$-anisotropic rather than isotropic.
\item \textbf{Uncertainty-guided adaptive budget.}
A monotone, bounded allocation rule that maps predicted trace
$\mathrm{tr}(\bm\Sigma)$ to sample count $K$ and polish budget $T$,
so difficult instances receive more evaluations without inflating the
budget on easy instances.
\item \textbf{Local guarantees.}
A suite of local results covering: the loss from trust-region
confinement (\cref{sec:trust_region}); gradient anti-concentration
under assumption~(A2); ordering preservation under depolarizing noise;
adaptive-allocation regret; spectral-norm generalization; and
finite-sample conformal coverage.
\item \textbf{Empirical evaluation.}
Ablation, cross-size transfer, family generalization, operating-point
sweep, multi-seed stability, and finite-shot simulation, all at
$p{=}2$ in the CPU-accessible regime.
\end{enumerate}

The paper is organized as follows.
\Cref{sec:background} collects notation and background.
\Cref{sec:related} situates the work relative to prior transfer,
uncertainty, and amortized-optimization literature.
\Cref{sec:method} describes the five algorithmic components.
\Cref{sec:theory} states and proves the local guarantees.
\Cref{sec:setup,sec:results} describe the experimental protocol and
empirical evaluation.
\Cref{sec:discussion} discusses implications and limitations.
\Cref{sec:conclusion} concludes.
Appendices contain proofs, implementation details, and supplementary
statistics.


\section{Problem Setting and Preliminaries}
\label{sec:background}

We work with unweighted MaxCut on $n$-vertex graphs and the
depth-$p$ \QAOA{}.

\subsection{MaxCut and the \QAOA{} circuit}

Let $\calG = (\calV,\calE)$ be an unweighted graph with
$|\calV| = n$ and $|\calE| = m$.  MaxCut asks for a bipartition
$z \in \{+1,-1\}^n$ maximizing
\begin{equation}
    C(z) = \sum_{(i,j) \in \calE} \tfrac{1}{2}(1 - z_i z_j).
    \label{eq:maxcut}
\end{equation}

The depth-$p$ \QAOA{} prepares the state
\begin{equation}
    \ket{\bm\gamma,\bm\beta}
    = e^{-i\beta_p B} e^{-i\gamma_p C} \cdots
    e^{-i\beta_1 B} e^{-i\gamma_1 C}
    \ket{+}^{\otimes n},
    \label{eq:qaoa_state}
\end{equation}
where $C = \sum_{(i,j)\in\calE}\frac{1}{2}(I - Z_i Z_j)$ is the
cost Hamiltonian and
$B = \sum_i X_i$~\cite{farhi2014}.  The $2p$ angles
$\bm\theta = (\bm\gamma,\bm\beta) \in \bbR^{2p}$ are chosen to
maximize
$F_\calG(\bm\theta) = \bra{\bm\theta}C\ket{\bm\theta}$.

Two quality measures appear below.  The analysis uses the normalized
expected objective $F_\calG(\bm\theta)/C_{\max}$.  The experiments
report the sampled approximation ratio
$r = C(z_{\mathrm{best}})/C_{\max}$, where $z_{\mathrm{best}}$ is the
best bitstring observed during optimization.

\subsection{Angle concentration and barren plateaus}
\label{sec:concentration_bg}

At fixed depth $p$ on bounded-degree graph families, optimal \QAOA{}
angles concentrate: the variance of the optimal $\bm\theta$ decreases
as $O(1/n)$~\cite{brandao2018,akshay2021}.  This makes cross-instance
parameter transfer possible~\cite{wurtz2021,galda2021}.

Variational circuits can also exhibit barren plateaus, where the
gradient variance at random parameters satisfies
\begin{equation}
    \mathrm{Var}_{\bm\theta}\bigl[\partial_j F(\bm\theta)\bigr]
    \le O\bigl(2^{-n}\bigr)
    \label{eq:bp}
\end{equation}
for deep or highly entangling
ans\"atze~\cite{mcclean2018,arrasmith2021,wang2021noise}.
Constant-depth \QAOA{} with a local cost operator avoids this
exponential suppression, but the landscape can still contain extended
flat regions.

\subsection{Spectral positional encodings}
\label{sec:spectral_bg}

The graph Laplacian $L = D - A$ has eigendecomposition
$L = U\Lambda U^\top$ with eigenvalues
$0 = \lambda_0 \le \lambda_1 \le \cdots \le \lambda_{n-1}$.
The eigenvectors $\{u_k\}$ define spectral positional
encodings~\cite{dwivedi2023,lim2023,kreuzer2021}. Sign ambiguity is
resolved by random sign flips during training~\cite{lim2023}.

\subsection{Graph Isomorphism Networks}

The Graph Isomorphism Network (\GIN)~\cite{xu2019} is as expressive
as the Weisfeiler--Leman isomorphism test~\cite{weisfeiler1968}.
Each layer updates
\begin{equation}
    h_v^{(l+1)} = \mathrm{MLP}^{(l)}\!\bigl(
    (1+\epsilon^{(l)})\, h_v^{(l)}
    + {\textstyle\sum_{u \in \calN(v)}} h_u^{(l)}
    \bigr),
    \label{eq:gin_update}
\end{equation}
with learnable $\epsilon^{(l)}$ and neighbor set
$\calN(v)$~\cite{hamilton2017,gilmer2017}.
A graph-level representation is obtained by permutation-invariant
pooling of the node embeddings.

\subsection{Wasserstein distance for Gaussians}
\label{sec:wasserstein_bg}

The 2-Wasserstein distance between probability measures $P$ and $Q$
on $\bbR^d$ is
\begin{equation}
    W_2(P,Q) = \Bigl(\inf_{\pi \in \Pi(P,Q)}
    \int \|\bm{x} - \bm{y}\|_2^2\, d\pi(\bm{x},\bm{y})\Bigr)^{1/2}.
    \label{eq:w2_def}
\end{equation}
For diagonal Gaussians this reduces to
\begin{equation}
    W_2^2 = \|\bm\mu_1 - \bm\mu_2\|_2^2
    + \|\bm\sigma_1 - \bm\sigma_2\|_2^2,
\end{equation}
which is closed form for diagonal covariance.


\section{Relation to Prior Work}
\label{sec:related}

\textbf{Parameter transfer for \QAOA{}.}
Angle concentration at fixed depth~\cite{brandao2018,akshay2021}
supports fixed schedules~\cite{wurtz2021}, interpolation rules~\cite{zhou2020},
cross-instance transfer~\cite{galda2021}, warm starts from related
circuits~\cite{egger2021,sack2021}, graph-conditioned predictors~\cite{jain2022,khairy2020},
Bayesian optimization~\cite{tibaldi2023}, and related search
strategies~\cite{shaydulin2023,streif2020,moussa2022,sauvage2024}.
These methods provide angles or search heuristics.  Here the predictor
produces a full distribution.  A point prediction can only initialize
search; a Gaussian prediction can also constrain search geometry and
allocate the per-instance budget.

\textbf{Learned combinatorial heuristics.}
In classical combinatorial optimization, learned heuristics are judged
mainly by wall-clock time to solution~\cite{bengio2021,cappart2023,khalil2017,dai2016,kool2019,schuetz2022}.
For variational quantum algorithms the dominant resource is often the
number of circuit evaluations.  Classical methods therefore target
throughput directly, whereas the relevant metric here is query count,
independent of processor speed.  Evaluation reduction is the primary
figure of merit, not wall-clock time.

\textbf{Uncertainty quantification.}
Standard tools include heteroscedastic regression, ensembles,
dropout-based approximations, and post hoc calibration~\cite{nix1994,kendall2017,lakshminarayanan2017,gal2016,kuleshov2018,abdar2021,hullermeier2021,psaros2023}.
Here uncertainty is used as an algorithmic quantity rather than only
as a diagnostic.  It defines the trust region and sets the evaluation
budget.  The Wasserstein regularizer and the contrastive loss are
introduced for this purpose.  A poorly calibrated variance head
produces a misaligned trust region and an incorrect budget, both of
which degrade optimization rather than merely corrupting a confidence
statement.

\textbf{Optimal transport and contrastive learning.}
The closed-form Wasserstein distance for
Gaussians~\cite{dowson1982} has been used in generative
modeling~\cite{arjovsky2017,frogner2015} and distributional
robustness~\cite{villani2009,cuturi2013,peyre2019}.  Here it
regularizes the predicted angle distribution.  Contrastive
objectives~\cite{chen2020simclr,oord2018,you2020,suresh2021,khosla2020}
structure the graph embedding by landscape similarity.  The two
regularizers act at different levels.  The contrastive loss acts on
the embedding before the Gaussian heads; it orders latent space by
optimization similarity.  The Wasserstein loss acts on the output
distribution; it couples the predicted covariance across instances
that share nearby angle optima.

\textbf{Amortized optimization.}
Amortized methods replace part or all of per-instance optimization by
a learned map~\cite{amos2023,chen2022,finn2017,shu2022}.  Existing
amortized \QAOA{} methods predict point initializations.  Here the
prediction is a Gaussian.  Learning does not replace optimization; it
restricts it.  This distinction affects the failure mode: an
amortized point predictor that transfers poorly leaves the optimizer
without guidance, whereas a Gaussian predictor that transfers poorly
enlarges the trust region and degrades toward the point-predictor
baseline rather than toward random initialization.

% ============================================================
\section{Graph-Conditioned Gaussian Prediction and Trust-Region Search}
\label{sec:method}
% ============================================================

A single forward pass of a \GIN{} produces a Gaussian
$(\bm\mu,\bm\Sigma)$ over \QAOA{} angles.  The mean sets the center,
the covariance sets the geometry, and the trace sets the budget.  The
method has five components (\cref{fig:algorithm}): a \GIN{} predictor,
Wasserstein-regularized training, contrastive metric learning,
trust-region-constrained optimization, and uncertainty-guided budget
allocation.

Each component serves a distinct algorithmic role.  The predictor maps
graph structure to an angle distribution.  The Wasserstein term aligns
nearby landscapes in distribution space.  The contrastive term orders
the representation by angle similarity before the output heads act.
The trust region converts predicted covariance into a geometric search
constraint.  The adaptive budget converts the same covariance into a
resource-allocation rule.  The method therefore uses one learned
object for initialization, confinement, and budget selection.  This is
the central structural point of the algorithm.  Prediction does not end
at the first iterate.  It persists through the full search procedure by
setting the center, the admissible anisotropy, and the amount of local
work invested on the instance.

\paragraph{Standing assumptions.}
The guarantees in \cref{sec:theory} are local. They assume:
(A1) $F_\calG$ is Lipschitz-continuous on the trust region
(Proposition~\ref{prop:lipschitz});
(A2) $\lambda_{\min}(-\nabla^2 F_\calG(\bm\mu)) \ge \kappa > 0$
(Theorem~\ref{thm:anticoncentration}); and
(A3) depolarizing noise with per-layer strength
$\varepsilon < 1$ (Theorem~\ref{thm:noise_robust}).
Assumption~(A2) is local.

\subsection{Gaussian predictor architecture}
\label{sec:architecture}

Given $\calG = (\calV,\calE)$, node features are formed by concatenating
normalized degree $d_v/n$ with $k$-dimensional Laplacian spectral
positional encodings~\cite{dwivedi2023,lim2023,kreuzer2021}, i.e., the
$k$ smallest nontrivial eigenvectors of $L = D - A$, yielding
$h_v^{(0)} \in \bbR^{k+1}$. Eigenvector sign ambiguity is handled by
random sign flips during training~\cite{lim2023}.

A three-layer \GIN{} with hidden dimension $d{=}64$ and
LayerNorm~\cite{ba2016} produces node embeddings. A graph-level
representation $\bm{g} \in \bbR^{2d}$ is obtained by concatenating mean
and max pooling:
\begin{equation}
    \bm{g} = \bigl[\tfrac{1}{n}\textstyle\sum_v h_v^{(L)} \;\big\|\;
    \max_v h_v^{(L)}\bigr].
    \label{eq:graph_repr}
\end{equation}
Two linear heads produce the Gaussian parameters:
\begin{align}
    \bm\mu &= W_\mu \bm{g} + b_\mu \in \bbR^{2p}, \label{eq:mu_head} \\
    \log\bm\sigma^2 &= \mathrm{clamp}(W_\sigma \bm{g} + b_\sigma,
    -5, 2). \label{eq:sigma_head}
\end{align}
The predictive distribution is
\begin{equation*}
    p(\bm\theta \mid \calG)
    = \Normal\bigl(\bm\mu(\calG),\,
    \mathrm{diag}(\bm\sigma^2(\calG))\bigr).
\end{equation*}

The clamp in \cref{eq:sigma_head} keeps
$\sigma_j \in [e^{-5/2}, e]$.  The covariance is diagonal.
The backbone is \GIN{} and the output distribution is diagonal
Gaussian.

The architecture couples graph representation and angle uncertainty in
the smallest form that still supports the algorithm.  Spectral
positional encodings supply global graph coordinates that persist
across graph size.  The \GIN{} backbone preserves local combinatorial
structure under message passing.  Mean pooling stabilizes family-level
statistics; max pooling preserves sharp structural signatures.  The
mean head sets a candidate center for optimization.  The variance head
sets the anisotropy of the search.  A point predictor would supply the
first object but not the second, and the trust-region mechanism would
collapse to unconstrained local search.  The diagonal Gaussian is not a
claim about full posterior structure.  It is the least complicated
object that still carries directional scale, and directional scale is
enough to define an ellipsoidal search set and an instance-dependent
budget rule.

\subsection{Wasserstein-regularized training}
\label{sec:training}

\textbf{Data.}
Training uses 240 graphs at $n{=}14$ (60 per family) and a validation
set of 80 graphs. Targets $\bm\theta^*_i$ are obtained by multi-restart
Nelder--Mead~\cite{nelder1965} on exact simulation.

\textbf{Phase 1.}
\begin{equation}
    \calL_1 = \frac{1}{N}\sum_{i=1}^{N}
    \|\bm\mu(\calG_i) - \bm\theta^*_i\|_2^2
    + \lambda_C \calL_{\mathrm{CL}}.
    \label{eq:loss_phase1}
\end{equation}

\textbf{Phase 2.}
\begin{equation}
    \calL_2 = \calL_{\mathrm{NLL}}
    + \lambda_W \calL_{\mathrm{W2}}
    + \lambda_C \calL_{\mathrm{CL}},
    \label{eq:loss_phase2}
\end{equation}

The negative log-likelihood is
\begin{equation}
    \calL_{\mathrm{NLL}}
    = \frac{1}{N}\sum_{i=1}^{N}\sum_{j=1}^{2p}
    \left[
    \frac{(\mu_j(\calG_i) - \theta^*_{ij})^2}
    {\sigma_j^2(\calG_i)}
    + \log\sigma_j^2(\calG_i)
    \right].
    \label{eq:loss_nll}
\end{equation}


\textbf{Wasserstein regularizer.}
For diagonal Gaussians,
\begin{equation*}
    W_2^2(P_i,P_j)
    = \|\bm\mu_i - \bm\mu_j\|_2^2
    + \|\bm\sigma_i - \bm\sigma_j\|_2^2.
\end{equation*}
Weights
$w_{ij} = \exp(-\|\bm\theta_i^* - \bm\theta_j^*\|_2^2 / 2\tau_W^2)$
encode angle-space similarity:
\begin{equation}
    \calL_{\mathrm{W2}} = \frac{1}{|\calB|^2}
    \sum_{i \ne j \in \calB} w_{ij}\, W_2^2(P_i,P_j).
    \label{eq:loss_w2}
\end{equation}
Phase~1 fits the mean.  Phase~2 fits the full Gaussian.

The two phases separate two tasks that otherwise interfere.  Phase~1
anchors the center of the distribution near the target angle vector.
Phase~2 assigns scale after the center is already stable.  The
negative log-likelihood alone acts pointwise in the training set.  The
Wasserstein term adds a geometric coupling across instances.  Two
graphs with nearby target angles then receive nearby predictive
distributions rather than unrelated variances.  That coupling matters
because the covariance is not a passive uncertainty label; it later
defines both the trust region and the budget.  Training therefore does
not merely fit labels.  It shapes a search geometry that will later be
used by the optimizer, so errors in scale become algorithmic errors
rather than cosmetic calibration errors.

\subsection{Contrastive metric learning}
\label{sec:contrastive}

A supervised contrastive loss structures the embedding $\bm{g}$. For
anchor $i$, define
\begin{equation*}
    \calP(i)
    = \bigl\{j \ne i :
    \|\bm\theta_i^* - \bm\theta_j^*\|_2 < \delta\bigr\}.
\end{equation*}
\begin{equation}
    \calL_{\mathrm{CL}} = -\frac{1}{|\calB|}\sum_i
    \frac{1}{|\calP(i)|}\sum_{j \in \calP(i)}
    \log\frac{e^{\mathrm{sim}(\bm{g}_i,\bm{g}_j)/\tau_C}}
    {\sum_{k \ne i} e^{\mathrm{sim}(\bm{g}_i,\bm{g}_k)/\tau_C}},
    \label{eq:loss_cl}
\end{equation}

This term acts before Gaussian prediction.  Its role is to give the
embedding a landscape order rather than a family label.  Nearby angle
optima pull graph representations together; distant optima push them
apart.  The prediction heads therefore start from a representation
whose local neighborhoods already reflect optimization similarity.  The
Wasserstein term then refines that structure at the level of the final
distribution.  In this sense, contrastive learning organizes the latent
space, and Wasserstein regularization organizes the output space.

\subsection{Trust-region \QAOA{} optimization}
\label{sec:trust_region}

\begin{definition}[Trust region]
\label{def:trust_region}
Given a predicted Gaussian $\Normal(\bm\mu,\bm\Sigma)$ for graph
$\calG$, with $\bm\Sigma$ symmetric positive definite, and a confidence
level $\alpha \in (0,1)$, the
\emph{trust region} is the Mahalanobis ellipsoid
\begin{equation}
    \calT_\alpha(\calG) = \bigl\{\bm\theta \in \bbR^{2p} :
    (\bm\theta - \bm\mu)^\top \bm\Sigma^{-1}
    (\bm\theta - \bm\mu) \le \chi^2_{2p}(\alpha)\bigr\},
    \label{eq:trust_region}
\end{equation}
where $\chi^2_{2p}(\alpha)$ is the $\alpha$-quantile of the
$\chi^2$ distribution with $2p$ degrees of freedom.
\end{definition}

Optimization is restricted to $\calT_\alpha$:
(a)~candidate angles
$\bm\theta_1,\ldots,\bm\theta_K \sim \Normal(\bm\mu,\bm\Sigma)$
are sampled from the corresponding truncated Gaussian restricted to
$\calT_\alpha$;
(b)~during Nelder--Mead refinement, any simplex vertex
$\bm\theta'$ exiting $\calT_\alpha$ is projected back:
\begin{equation}
    \mathrm{proj}_{\calT}(\bm\theta') = \bm\mu +
    \min\!\Bigl(1,\;
    \frac{\sqrt{\chi^2_{2p}(\alpha)}}
    {\|\bm\Sigma^{-1/2}(\bm\theta' - \bm\mu)\|_2}
    \Bigr)(\bm\theta' - \bm\mu).
    \label{eq:projection}
\end{equation}

The trust region is the algorithmic core.  It turns predicted
covariance into an admissible search set.  A warm start uses the mean
once and then discards the predictor.  The present scheme keeps the
prediction active through the full refinement stage.  Every iterate is
measured in the Mahalanobis geometry induced by $\bm\Sigma$.  Small
variance directions receive tight control; large variance directions
remain available for exploration.  The search therefore inherits the
anisotropy of the learned angle distribution rather than a fixed,
hand-chosen radius.  This distinction is the source of the query
savings.  The algorithm does not spend evaluations discovering a local
metric from scratch around each new graph.  It begins with a metric
predicted from related solved instances and restricts search to that
geometry from the first query onward.

\begin{figure}[t]
    \centering
    \includegraphics[width=\columnwidth]{figures/fig_landscape_trust.pdf}
    \caption{Mechanism of the learned trust region on an $n{=}8$
        Erd\H{o}s--R\'enyi instance at $p{=}2$ (layer-1 slice).
        Panel (a) shows the local \QAOA{} landscape, panel (b) the
        learned Mahalanobis ellipsoids at $95\%$ and $68\%$ coverage,
        and panel (c) the contrast between unconstrained random
        initializations (gray) and samples restricted to the learned
        region (red).  The figure emphasizes the geometric content of
        the predictor: it supplies an anisotropic local metric that
        removes evaluations spent exploring directions already judged
        implausible.
    \label{fig:landscape_trust}}
\end{figure}

\subsection{Uncertainty-guided adaptive refinement}
\label{sec:adaptive}

Define
\begin{equation*}
    U(\calG) = \mathrm{tr}(\bm\Sigma(\calG))/2p,
    \quad
    z(\calG) = \frac{U(\calG) - U_{\mathrm{med}}}{U_{\mathrm{iqr}}}.
\end{equation*}

Resource allocation:
\begin{align}
    K(\calG) &= \mathrm{clamp}(\lfloor 1 + 4\sigma(z)\rfloor,1,5),
    \label{eq:K_adaptive} \\
    T(\calG) &= \mathrm{clamp}(\lfloor T_{\mathrm{base}}
    (0.5 + z_+)\rfloor, 5, 2T_{\mathrm{base}}),
    \label{eq:T_adaptive}
\end{align}

Higher uncertainty increases both the number of samples and the
number of refinement steps.  The mapping is monotone and bounded
(\cref{fig:budget_alloc}).

This allocation rule gives the variance head a second algorithmic use.
The same signal that enlarges the search region also increases query
budget.  Easy instances therefore receive a narrow region and short
refinement; difficult instances receive a broader region and a larger
budget.  A uniform budget would ignore the instance-dependent content
of the predictor and spend queries on graphs that already admit a
stable center.  The allocation rule does not solve a separate control
problem.  It enforces structural consistency between predicted geometric
spread and computational effort: large predicted spread enlarges both
the search set and the work invested inside it.

\begin{figure}[t]
\noindent\fbox{\parbox{0.95\columnwidth}{\raggedright%
\textbf{Algorithm 1: \UQQAOA{} inference}\\[4pt]%
\textbf{Input:} graph $\calG$, trained model
$f_\theta : \calG \mapsto (\bm\mu,\bm\Sigma)$,
budget $T_{\mathrm{base}}$, confidence level $\alpha$.\\[3pt]
\textbf{Output:} approximate MaxCut bitstring $z^*$.\\[3pt]
1. Predict $(\bm\mu, \bm\Sigma) \leftarrow f_\theta(\calG)$.\\
2. Compute $U \leftarrow \mathrm{tr}(\bm\Sigma)/2p$;\;
   $z \leftarrow (U - U_{\mathrm{med}})/U_{\mathrm{iqr}}$.\\
3. Set sample count
   $K \leftarrow \mathrm{clamp}(\lfloor 1 + 4\sigma(z)\rfloor,1,5)$.\\
4. Set polish budget
   $T \leftarrow \mathrm{clamp}(\lfloor T_{\mathrm{base}}(0.5+z_+)
   \rfloor, 5, 2T_{\mathrm{base}})$.\\
5. Construct trust region
   $\calT_\alpha \leftarrow \{\bm\theta :
   \|\bm\Sigma^{-1/2}(\bm\theta - \bm\mu)\|_2^2
   \le \chi^2_{2p}(\alpha)\}$.\\
6. Sample $\bm\theta_1,\ldots,\bm\theta_K$ from the truncated
    Gaussian $\Normal(\bm\mu,\bm\Sigma)$ restricted to $\calT_\alpha$.\\
7. Evaluate \QAOA{} circuit: $\bm\theta^* \leftarrow
   \arg\max_k F_\calG(\bm\theta_k)$.\\
8. Trust-region polish: run $T$ iterations of Nelder--Mead
   from $\bm\theta^*$, projecting each iterate via
   \cref{eq:projection}.\\
9. Return best bitstring $z^*$ from final \QAOA{} measurement.
}}
\caption{\UQQAOA{} inference.  The diagram makes explicit that a
single predicted Gaussian is reused three times at test time: its
mean initializes search, its covariance defines the trust region, and
its trace sets the sampling and refinement budgets.  Steps~3--4
convert uncertainty into compute allocation, while Steps~5--8 execute
trust-region-constrained optimization.
\label{fig:algorithm}}
\end{figure}

\begin{remark}[Conformal trust-region radius]
\label{rem:conformal_tr}
At inference time, the $\chi^2_{2p}(\alpha)$ radius in Step~5 can
be replaced by the conformal quantile $\hat{q}_{1-\alpha}$
(\cref{prop:conformal}), yielding a \emph{conformalized} trust
region
\begin{equation*}
    \calT_{\hat{q}_{1-\alpha}}
    = \bigl\{\bm\theta :
    (\bm\theta - \bm\mu)^\top \bm\Sigma^{-1}(\bm\theta - \bm\mu)
    \le \hat{q}_{1-\alpha}\bigr\}
\end{equation*}
with the distribution-free, finite-sample guarantee
$\bbP[\bm\theta^* \in \calT_{\hat{q}_{1-\alpha}}] \ge 1 - \alpha$.
For our validation set ($\hat{q}_{0.90} = 10.8$ vs.\
$\chi^2_4(0.90) = 7.779$), the radius increases by a
factor of $(10.8/7.779)^{1/2} \approx 1.178$.
\end{remark}

% ============================================================
\section{Guarantees for Learned Trust-Region Search}
\label{sec:theory}
% ============================================================

Each result below is local and conditional on the stated assumptions.
There is no claim of global optimality.  The bounds quantify the loss
incurred by restricting search to the predicted ellipsoid and by
allocating budget according to predicted uncertainty.

Throughout, $\calG \sim \calD$ denotes a held-out graph except where
otherwise stated. The predictor is
$f_\theta(\calG) = (\bm\mu(\calG),\bm\Sigma(\calG))$.
The reference point $\bm\theta^*(\calG)$ denotes the local optimum
produced by the multi-restart Nelder--Mead procedure used for target
generation, or the optimum specified in the corresponding statement.
The trust region $\calT_\alpha(\calG)$ is defined in
\cref{eq:trust_region}. For a radius not tied to a
$\chi^2_{2p}(\alpha)$ quantile, we write $\calT_q(\calG)$.

Expectations and probabilities are taken with respect to the randomness
explicitly indicated: draws of $\calG$ from $\calD$, samples from
$\Normal(\bm\mu(\calG),\bm\Sigma(\calG))$, or uniform samples on
$\calT_\alpha(\calG)$. Constants such as $L_{\calG,K}$, $L_\calG$,
$\kappa$, $\Lambda$, and $M_3$ are deterministic once $\calG$ is fixed.

\subsection{Landscape properties}

\begin{proposition}[Local Lipschitz regularity of the \QAOA{} objective]
\label{prop:lipschitz}
For a connected graph $\calG$ with $m$ edges, the depth-$p$ \QAOA{}
objective $F_\calG : \bbR^{2p} \to \bbR$ is real analytic. In
particular, on any bounded set $K \subset \bbR^{2p}$,
\begin{equation}
    L_{\calG,K} := \sup_{\bm\theta \in K}
    \|\nabla F_\calG(\bm\theta)\|_2 < \infty.
    \label{eq:local_lipschitz_def}
\end{equation}
A conservative global bound is $L_\calG \le 2m\sqrt{2p}$.
\end{proposition}

\begin{proof}
The \QAOA{} unitary is
$U(\bm\theta) = \prod_{\ell=1}^{p} e^{-i\beta_\ell B} e^{-i\gamma_\ell C}$,
where $B = \sum_v X_v$ and $C = \sum_{(i,j)\in\calE}\frac{1}{2}(I-Z_iZ_j)$.
Each factor is a matrix exponential of a bounded Hermitian operator, so
$U(\bm\theta)$ is an entire function of each angle.  The objective
$F_\calG(\bm\theta) = \bra{+}U^\dagger(\bm\theta)\,C\,U(\bm\theta)\ket{+}$
is therefore a real-analytic function on $\bbR^{2p}$, and
$\nabla F_\calG$ is continuous everywhere.  On any compact set $K$ the
continuous function $\|\nabla F_\calG\|_2$ attains its supremum, so
$L_{\calG,K} < \infty$.  The mean-value theorem then gives
$|F_\calG(\bm\theta) - F_\calG(\bm\theta')| \le L_{\calG,K}
\|\bm\theta - \bm\theta'\|_2$ for all
$\bm\theta,\bm\theta' \in K$.
For the global bound, note that each partial derivative of $F_\calG$
is an expectation value of a commutator.  For the phase-separation
angle $\gamma_\ell$,
$\partial_{\gamma_\ell} F_\calG
= -i\bra{\bm\theta}[C_\ell, U_\ell^\dagger C U_\ell]\ket{\bm\theta}$
where $C_\ell$ is the cost operator acting in layer~$\ell$.  Since $C$
has operator norm $\|C\|_{\mathrm{op}} \le m$ (each edge contributes at
most~1), the commutator has norm at most $2m$, giving
$|\partial_{\gamma_\ell} F_\calG| \le m$.  Similarly,
$|\partial_{\beta_\ell} F_\calG| \le n$ since $\|B\|_{\mathrm{op}} = n$.
The full gradient satisfies
$\|\nabla F_\calG\|_2^2
= \sum_{\ell=1}^p [(\partial_{\gamma_\ell} F)^2
+ (\partial_{\beta_\ell} F)^2]
\le p(m^2 + n^2)$.
For connected graphs, $m \ge n-1$ so $m^2 + n^2 \le 2m^2 + 2m + 1
\le 4m^2$ for $m \ge 2$.  Hence
$\|\nabla F_\calG\|_2 \le 2m\sqrt{p} \le 2m\sqrt{2p}$,
which establishes the global bound.
\end{proof}

\subsection{Trust-region guarantees}

\begin{lemma}[Validity of the Mahalanobis projection]
\label{lem:projection_valid}
Let $\bm\mu \in \bbR^{2p}$,
$\bm\Sigma \succ 0$, and $q > 0$.  For any
$\bm\theta' \in \bbR^{2p}$, the map
\begin{equation*}
    \mathrm{proj}_{\calT}(\bm\theta') = \bm\mu +
    \min\!\Bigl(1,\;
    \frac{\sqrt{q}}
    {\|\bm\Sigma^{-1/2}(\bm\theta' - \bm\mu)\|_2}
    \Bigr)(\bm\theta' - \bm\mu)
\end{equation*}
satisfies:
\begin{enumerate}[label=(\roman*)]
\item $\mathrm{proj}_{\calT}(\bm\theta') \in \calT_q$ for every
$\bm\theta' \in \bbR^{2p}$.
\item If $\bm\theta' \in \calT_q$, then
$\mathrm{proj}_{\calT}(\bm\theta') = \bm\theta'$ (idempotency).
\item $\mathrm{proj}_{\calT}$ is the nearest point in $\calT_q$ under
the Mahalanobis norm
$\|\cdot\|_{\bm\Sigma^{-1}}$.
\end{enumerate}
\end{lemma}

\begin{proof}
Write $\bm d = \bm\theta' - \bm\mu$ and
$s = \|\bm\Sigma^{-1/2}\bm d\|_2$.  The projected point is
$\bm\theta_{\mathrm{pr}} = \bm\mu + \min(1, \sqrt{q}/s)\,\bm d$.

\emph{(i)} If $s \le \sqrt{q}$, then $\min(1,\sqrt{q}/s) = 1$ and
$\bm\theta_{\mathrm{pr}} = \bm\theta'$, which satisfies the constraint
$s^2 \le q$.  If $s > \sqrt{q}$,
\begin{equation*}
    \|\bm\Sigma^{-1/2}(\bm\theta_{\mathrm{pr}} - \bm\mu)\|_2
    = \frac{\sqrt{q}}{s}\,s = \sqrt{q},
\end{equation*}
so $\bm\theta_{\mathrm{pr}}$ lies on the boundary of $\calT_q$.

\emph{(ii)} Immediate from the first case above.

\emph{(iii)} For $s > \sqrt{q}$, the projection scales $\bm d$ to the
boundary along the direction $\bm d / s$.  The Mahalanobis ball
$\calT_q$ is convex, and radial scaling toward the center is the
Euclidean projection in the whitened coordinates
$\bm u = \bm\Sigma^{-1/2}(\bm\theta - \bm\mu)$, which is the
nearest-point projection onto the ball $\|\bm u\|_2 \le \sqrt{q}$.
\end{proof}

\begin{theorem}[Lower bound on the expected objective within the trust region]
\label{thm:landscape_lb}
Let $\calT_\alpha(\calG)$ be defined by \cref{eq:trust_region}. Suppose
$F_\calG(\bm\mu) \ge r^* C_{\max}$. Let $\Theta$ be supported on
$\calT_\alpha(\calG)$. Then
\begin{equation}
    \bbE\bigl[F_\calG(\Theta)\bigr]
    \ge r^* C_{\max}
    - L_\calG \sqrt{\chi^2_{2p}(\alpha)\cdot
    \|\bm\Sigma\|_{\mathrm{op}}},
    \label{eq:landscape_lb}
\end{equation}
\end{theorem}
\begin{proof}
For $\bm\theta \in \calT_\alpha$, one has
\begin{equation*}
    \|\bm\theta - \bm\mu\|_2^2
    \le \|\bm\Sigma\|_{\mathrm{op}}\cdot
    (\bm\theta - \bm\mu)^\top \bm\Sigma^{-1}(\bm\theta - \bm\mu)
    \le \|\bm\Sigma\|_{\mathrm{op}}\cdot \chi^2_{2p}(\alpha).
\end{equation*}
Hence
\begin{equation*}
    \|\bm\theta - \bm\mu\|_2
    \le \sqrt{\chi^2_{2p}(\alpha)\|\bm\Sigma\|_{\mathrm{op}}}.
\end{equation*}
By Proposition~\ref{prop:lipschitz},
\begin{equation*}
    |F_\calG(\bm\theta) - F_\calG(\bm\mu)|
    \le L_\calG \sqrt{\chi^2_{2p}(\alpha)\|\bm\Sigma\|_{\mathrm{op}}}.
\end{equation*}
Applying the assumption $F_\calG(\bm\mu) \ge r^* C_{\max}$ and taking
expectation yields the claim.
\end{proof}

\begin{corollary}[Expected objective under variance confinement]
\label{cor:quality_confinement}
Under the conditions of Theorem~\ref{thm:landscape_lb}, if
$\sigma_{\max}^2 \le (r^* C_{\max})^2 /
(4 L_\calG^2 \chi^2_{2p}(\alpha))$, then
$\bbE[F_\calG(\Theta)] \ge \tfrac{1}{2} r^* C_{\max}$.
\end{corollary}

\begin{proof}
Set $\epsilon = \tfrac{1}{2}r^* C_{\max}$ in
Theorem~\ref{thm:landscape_lb}. This gives the sufficient condition
$\sigma_{\max}^2 \le (r^* C_{\max})^2 /
(4 L_\calG^2 \chi^2_{2p}(\alpha))$.
The conclusion follows immediately.
\end{proof}

\begin{proposition}[Second-order local improvement within the trust region]
\label{prop:second_order}
Assume $F_\calG$ is twice continuously differentiable near $\bm\mu$
with $\lambda_{\min}(-\nabla^2 F_\calG(\bm\mu)) \ge \kappa > 0$
(i.e., $\bm\mu$ is a strict local maximizer).  If the trust region
satisfies
$\sqrt{\chi^2_{2p}(\alpha)\,\|\bm\Sigma\|_{\mathrm{op}}}
\le \kappa / M_3$,
where $M_3$ bounds the third derivatives on $\calT_\alpha$, then for
every $\bm\theta \in \calT_\alpha$,
\begin{equation}
    \begin{aligned}
    F_\calG(\bm\theta)
    &\ge F_\calG(\bm\mu)
    - \tfrac{1}{2}\Lambda\,
    \chi^2_{2p}(\alpha)\,\|\bm\Sigma\|_{\mathrm{op}} \\
    &\quad - \tfrac{1}{6}M_3\,
    \bigl[\chi^2_{2p}(\alpha)\,
    \|\bm\Sigma\|_{\mathrm{op}}\bigr]^{3/2},
    \end{aligned}
    \label{eq:second_order_lb}
\end{equation}
where $\Lambda = \|\nabla^2 F_\calG(\bm\mu)\|_{\mathrm{op}}$.  For a small trust region the leading penalty is quadratic in
$\sigma_{\max}$ rather than linear.
\end{proposition}

\begin{proof}
By Taylor expansion around $\bm\mu$ (a critical point, so
$\nabla F_\calG(\bm\mu) = 0$ for a local maximizer):
\begin{equation*}
    F_\calG(\bm\theta)
    = F_\calG(\bm\mu)
    + \tfrac{1}{2}\Delta^\top \nabla^2 F_\calG(\bm\mu)\,\Delta
    + O(\|\Delta\|_2^3),
\end{equation*}
where $\Delta = \bm\theta - \bm\mu$.  Since $\bm\mu$ is a local
maximizer, $\nabla^2 F_\calG(\bm\mu) \preceq 0$, so
$\Delta^\top \nabla^2 F_\calG(\bm\mu)\,\Delta
\ge -\Lambda\|\Delta\|_2^2$.  The cubic remainder is bounded by
$\tfrac{1}{6}M_3(2p)^{3/2}\|\Delta\|_2^3
\le \tfrac{1}{6}M_3\|\Delta\|_2^3$ (absorbing dimension factors into
$M_3$).  For $\bm\theta \in \calT_\alpha$,
$\|\Delta\|_2^2 \le \chi^2_{2p}(\alpha)\,\|\bm\Sigma\|_{\mathrm{op}}$
(as in Theorem~\ref{thm:landscape_lb}).  Substituting:
\begin{align*}
    F_\calG(\bm\theta)
    &\ge F_\calG(\bm\mu)
    - \tfrac{1}{2}\Lambda\,\|\Delta\|_2^2
    - \tfrac{1}{6}M_3\,\|\Delta\|_2^3 \\
    &\ge F_\calG(\bm\mu)
    - \tfrac{1}{2}\Lambda\,\chi^2_{2p}(\alpha)\,
    \|\bm\Sigma\|_{\mathrm{op}} \\
    &\quad\; - \tfrac{1}{6}M_3\,
    [\chi^2_{2p}(\alpha)\,\|\bm\Sigma\|_{\mathrm{op}}]^{3/2}.
    \qedhere
\end{align*}
\end{proof}


\subsection{Gradient anti-concentration}

\begin{theorem}[Local gradient anti-concentration in the trust region]
\label{thm:anticoncentration}
Assume $F_\calG$ is three times continuously differentiable near
$\calT_\alpha(\calG)$,
$\lambda_{\min}(-\nabla^2 F_\calG(\bm\mu)) \ge \kappa > 0$,
and third derivatives are bounded by $M_3$.
Then for $\bm\theta \sim \mathrm{Unif}(\calT_\alpha)$,
\begin{equation}
    \begin{aligned}
    \mathrm{Var}_{\bm\theta \sim \mathrm{Unif}(\calT_\alpha)}
    \bigl[\partial_j F_\calG(\bm\theta)\bigr]
    &\ge \frac{\kappa^2\, \sigma_{\min}^2\, \chi^2_{2p}(\alpha)}
    {2p+2} \\
    &\quad - C_{p,\alpha}\bigl(\Lambda M_3\sigma_{\max}^3
    + M_3^2\sigma_{\max}^4\bigr),
    \end{aligned}
    \label{eq:anticoncentration}
\end{equation}
where $\Lambda := \|\nabla^2 F_\calG(\bm\mu)\|_{\mathrm{op}}$,
$\sigma_{\min} = \min_j \sigma_j$, and
$\sigma_{\max} = \max_j \sigma_j$.
\end{theorem}

\begin{proof}
See \cref{app:proof_anticonc}.
\end{proof}

\subsection{Algorithmic guarantees}

\begin{proposition}[Feasibility and monotonic best-so-far value under projection]
\label{prop:nondeg}
Projected optimization within $\calT_\alpha$ produces feasible iterates
and a nondecreasing best-so-far objective. The final iterate dominates
the initial sampled candidates. Let $\{\bm\theta_t\}_{t=0}^T$ be the
feasible iterates generated by projected Nelder--Mead within
$\calT_\alpha$, define $b_t := \max_{0 \le s \le t} F_\calG(\bm\theta_s)$,
and let $\widehat{\bm\theta}_t$ denote any feasible iterate attaining
$b_t$. Then:
\begin{enumerate}[label=(\roman*)]
\item $\bm\theta_t \in \calT_\alpha$ for all $t$.
\item $b_t$ is non-decreasing in $t$.
\item $F_\calG(\widehat{\bm\theta}_T) \ge \max_k F_\calG(\bm\theta_k)$
where $\bm\theta_1,\ldots,\bm\theta_K$ are the initial feasible samples
used to seed the search.
\end{enumerate}
\end{proposition}

\begin{proof}
(i)~The projection operator~\cref{eq:projection} maps every point
to $\calT_\alpha$: if $\bm\theta' \in \calT_\alpha$, the projection
is the identity; otherwise, it rescales
$\bm\theta' - \bm\mu$ to the boundary of $\calT_\alpha$.
(ii)~Monotonicity is immediate from the definition
$b_t = \max_{0 \le s \le t} F_\calG(\bm\theta_s)$: the index set on
which the maximum is taken only enlarges with~$t$.
(iii)~Because the initial sampled points are among the feasible points
seen by the algorithm before refinement begins,
$b_T \ge \max_{1 \le k \le K} F_\calG(\bm\theta_k)$. Since
$\widehat{\bm\theta}_T$ attains $b_T$, the claim follows.
\end{proof}

\begin{corollary}[Output guarantee under trust-region capture]
\label{cor:output_guarantee}
Let $\bm\theta_{\mathrm{out}}$ denote the best feasible angle vector returned by
Algorithm~\ref{fig:algorithm}, and let $z_{\mathrm{out}}$ be the final
bitstring sampled from the \QAOA{} circuit at
$\bm\theta_{\mathrm{out}}$. Fix a trust region $\calT_q(\calG)$ and
assume the reference local optimum $\bm\theta^*(\calG)$ belongs to that
region. Then
\begin{equation}
    F_\calG(\bm\theta_{\mathrm{out}})
    \ge F_\calG\bigl(\bm\theta^*(\calG)\bigr)
    - 2L_\calG\sqrt{q\,\|\bm\Sigma(\calG)\|_{\mathrm{op}}},
    \label{eq:output_angle_guarantee}
\end{equation}
and therefore
\begin{equation}
    \begin{aligned}
    \bbE\bigl[C(z_{\mathrm{out}})
    \mid \bm\theta^*(\calG) \in \calT_q(\calG)\bigr]
    &\ge F_\calG\bigl(\bm\theta^*(\calG)\bigr) \\
    &\quad - 2L_\calG\sqrt{q\,\|\bm\Sigma(\calG)\|_{\mathrm{op}}}.
    \end{aligned}
    \label{eq:output_bitstring_guarantee}
\end{equation}
If, in addition, the trust region has capture probability
$p_{\mathrm{cap}}(\calG)$, then since MaxCut costs are nonnegative,
\begin{equation}
    \bbE\bigl[C(z_{\mathrm{out}})\bigr]
    \ge p_{\mathrm{cap}}(\calG)
    \Bigl[F_\calG\bigl(\bm\theta^*(\calG)\bigr)
    - 2L_\calG\sqrt{q\,\|\bm\Sigma(\calG)\|_{\mathrm{op}}}\Bigr].
    \label{eq:output_unconditional}
\end{equation}
For the algorithmic $\chi^2$ region one may take
$p_{\mathrm{cap}} \ge \alpha-\eta$ from
Proposition~\ref{prop:trust_region_adequacy}; for the conformalized
region one has $p_{\mathrm{cap}} \ge 1-\alpha$.
\end{corollary}

\begin{proof}
Under the capture event, every sampled candidate and every projected
iterate lies in $\calT_q(\calG)$. Proposition~\ref{prop:trust_region_adequacy}
therefore gives the pointwise bound
\begin{equation*}
    F_\calG(\bm\theta)
    \ge F_\calG\bigl(\bm\theta^*(\calG)\bigr)
    - 2L_\calG\sqrt{q\,\|\bm\Sigma(\calG)\|_{\mathrm{op}}}
    \qquad \forall\,\bm\theta \in \calT_q(\calG).
\end{equation*}
Since Algorithm~\ref{fig:algorithm} returns an angle vector inside that
region, \cref{eq:output_angle_guarantee} follows. By definition of the
\QAOA{} objective,
$\bbE[C(z_{\mathrm{out}})\mid \bm\theta_{\mathrm{out}}]
= F_\calG(\bm\theta_{\mathrm{out}})$,
which gives \cref{eq:output_bitstring_guarantee}. The unconditional
bound \cref{eq:output_unconditional} follows by conditioning on the
capture event and using $C(z_{\mathrm{out}}) \ge 0$ for MaxCut.
\end{proof}

\begin{corollary}[Dominance over shared-center point initialization]
\label{cor:shared_center_dominance}
Assume the setting of Corollary~\ref{cor:output_guarantee}.  Suppose
the initial candidate set contains the predicted mean $\bm\mu(\calG)$.
Then the returned angle vector satisfies
\begin{equation}
    F_\calG(\bm\theta_{\mathrm{out}}) \ge F_\calG(\bm\mu(\calG)).
    \label{eq:shared_center_monotone}
\end{equation}
If, in addition, $\bm\theta^*(\calG) \in \calT_q(\calG)$, then
\begin{equation}
    F_\calG(\bm\theta_{\mathrm{out}})
    \ge F_\calG\bigl(\bm\theta^*(\calG)\bigr)
    - 2L_\calG\sqrt{q\,\|\bm\Sigma(\calG)\|_{\mathrm{op}}}.
    \label{eq:shared_center_target_gap}
\end{equation}
Thus, for a method that shares the same predicted center but does not
use the covariance to constrain search, trust-region refinement cannot
underperform the shared-center initialization and inherits the same
capture-conditioned target-gap bound.
\end{corollary}

\begin{proof}
Because the initial candidate set contains $\bm\mu(\calG)$ and
Proposition~\ref{prop:nondeg} states that the best-so-far objective is
nondecreasing over all feasible iterates, the returned point satisfies
\cref{eq:shared_center_monotone}.  The second claim is exactly the
output guarantee from Corollary~\ref{cor:output_guarantee}.
\end{proof}


\subsection{Sampling guarantees}

\begin{theorem}[Best-of-$K$ guarantee]
\label{thm:bestofk}
Let $\bm\theta_1,\ldots,\bm\theta_K \stackrel{\text{i.i.d.}}{\sim}
\Normal(\bm\mu,\bm\Sigma)$. Then
\begin{equation}
    \bbE\bigl[\max_k F_\calG(\bm\theta_k)\bigr]
    \ge F_\calG(\bm\mu) - \frac{L_\calG\sqrt{\mathrm{tr}(\bm\Sigma)}}
    {\sqrt{K}}.
    \label{eq:bestofk}
\end{equation}
\end{theorem}

\begin{proof}
See \cref{app:proof_bestofk}.
\end{proof}

\begin{lemma}[Trust-region volume ratio]
\label{lem:dim_red}
For a diagonal covariance $\bm\Sigma = \mathrm{diag}(\sigma_1^2,
\ldots,\sigma_{2p}^2)$ with confidence $\alpha{=}0.95$ and $p{=}2$,
the trust-region volume relative to the unconstrained search
hypercube $[-\pi,\pi]^{2p}$ is
\begin{equation}
    \frac{\mathrm{Vol}(\calT_{0.95})}{\mathrm{Vol}([-\pi,\pi]^{2p})}
    = \frac{\pi^p [\chi^2_{4}(0.95)]^p}
    {\Gamma(p{+}1)\,(2\pi)^{2p}}
    \prod_{j=1}^{4}\sigma_j.
    \label{eq:vol_ratio}
\end{equation}
If one uses the isotropic proxy $\sigma_j \equiv \bar\sigma = 0.15$,
then \cref{eq:vol_ratio} gives
$0.2850\,\bar\sigma^4 \approx 1.44 \times 10^{-4}$.
\end{lemma}

\begin{proof}
The trust region is a 4-dimensional ellipsoid with half-lengths
$\sigma_j \sqrt{\chi^2_4(0.95)}$.
Its volume is $\frac{\pi^2}{2}[\chi^2_4(0.95)]^2
\prod_j \sigma_j$,
and the hypercube volume is $(2\pi)^4$.
The ratio follows by direct computation with $\chi^2_4(0.95)
\approx 9.49$.
\end{proof}

\subsection{Generalization and calibration}


\begin{proposition}[Wasserstein continuity in a WL-compatible pseudometric]
\label{prop:wasserstein_lip}
If the \GIN{} encoder $\phi$ is $L_{\mathrm{enc}}$-Lipschitz
with respect to a WL-compatible graph pseudometric $d_{\mathrm{WL}}$
and the prediction heads are $L_{\mathrm{head}}$-Lipschitz, then:
\begin{equation}
    W_2\bigl(f_\theta(\calG),f_\theta(\calG')\bigr)
    \le \sqrt{2}\, L_{\mathrm{head}} L_{\mathrm{enc}}\,
    d_{\mathrm{WL}}(\calG,\calG'),
    \label{eq:wlip}
\end{equation}
where $d_{\mathrm{WL}}(\calG,\calG')=0$ is allowed for
non-isomorphic graphs that are indistinguishable by the relevant WL
statistics.
\end{proposition}

\begin{proof}
See \cref{app:proof_wlip}.
\end{proof}

\begin{theorem}[Adaptive-allocation guarantee under calibration and monotone difficulty]
\label{thm:regret}
Assume approximate calibration and the regret envelope
\cref{eq:monotone_difficulty_model}. Suppose the predicted uncertainty
satisfies
$\eta$-approximate calibration:
\begin{equation}
    \bigl|\bbP[\|\bm\theta^* - \bm\mu\|_{\bm\Sigma^{-1}} \le r]
    - F_{\chi^2_{2p}}(r^2)\bigr| \le \eta
    \quad \forall\, r > 0,
    \label{eq:eta_cal}
\end{equation}
where $F_{\chi^2_{2p}}$ is the $\chi^2_{2p}$ CDF\@.
Assume the residual optimization regret satisfies the instance-wise
envelope
\begin{equation}
    R_\pi(\calG) \le c_0 + c_1\frac{\sqrt{U(\calG)+\eta}}
    {\sqrt{K_\pi(\calG)}}
    \label{eq:monotone_difficulty_model}
\end{equation}
where $U(\calG)=\mathrm{tr}(\bm\Sigma(\calG))/2p$ and
$c_0,c_1 \ge 0$ are independent of $\pi$. Define population regret by
\begin{equation}
    \mathrm{Regret}(\pi) := \bbE_{\calG\sim\calD}[R_\pi(\calG)].
    \label{eq:policy_regret_def}
\end{equation}
Let $\bar K := \bbE_\calD[K_{\mathrm{adapt}}(\calG)]$, and let
$\pi_{\mathrm{unif}}$ denote the mean-budget-matched uniform comparator
with $K_{\mathrm{unif}}(\calG) \equiv \bar K$.
Assume further that the adaptive allocation rule satisfies the
envelope-dominance condition
\begin{equation}
    \bbE_\calD\!\left[
    \frac{\sqrt{U(\calG)+\eta}}{\sqrt{K_{\mathrm{adapt}}(\calG)}}
    \right]
    \le \frac{1}{\sqrt{\bar K}}
    \bbE_\calD\!\left[\sqrt{U(\calG)+\eta}\right].
    \label{eq:allocation_dominance}
\end{equation}
Then the adaptive allocation strategy $\pi_{\mathrm{adapt}}$
(Eqs.~\ref{eq:K_adaptive}--\ref{eq:T_adaptive}) satisfies
\begin{align}
    \mathrm{Regret}(\pi_{\mathrm{adapt}})
    &\le c_0 + c_1\,\bbE_\calD\!\left[
    \frac{\sqrt{U(\calG)+\eta}}{\sqrt{K_{\mathrm{adapt}}(\calG)}}
    \right] \nonumber \\
    &\le c_0 + \frac{c_1}{\sqrt{\bar K}}
    \bbE_\calD\!\left[\sqrt{U(\calG)+\eta}\right],
    \label{eq:regret} \\
    \bbE[B_{\pi_{\mathrm{adapt}}}]
    &= \alpha\,\bbE[B_{\pi_{\mathrm{unif}}}],
    \label{eq:budget_savings}
\end{align}
where $B_\pi$ is the total evaluation budget under policy~$\pi$ and
$\alpha = \bbE_\calD[B_{\pi_{\mathrm{adapt}}}] /
\bbE_\calD[B_{\pi_{\mathrm{unif}}}]$ depends on the deployment
distribution of predicted uncertainties. Thus $\alpha<1$ is not a
consequence of calibration alone; it depends on the deployment
distribution together with the implemented allocation rule.
\end{theorem}

\begin{proof}
See \cref{app:proof_regret}.
\end{proof}

\begin{proposition}[Monotonicity of the implemented allocation rule]
\label{prop:allocation_support}
Let
$z(\calG) = (U(\calG)-U_{\mathrm{med}})/U_{\mathrm{iqr}}$, and define
$K_{\mathrm{adapt}}$ and $T_{\mathrm{adapt}}$ by
\cref{eq:K_adaptive,eq:T_adaptive}. Then both allocations are
nondecreasing functions of $U$, and one has the explicit thresholds
\begin{equation}
    K_{\mathrm{adapt}}(\calG) =
    \begin{cases}
        1, & z(\calG) < -\log 3, \\
        2, & -\log 3 \le z(\calG) < 0, \\
        3, & 0 \le z(\calG) < \log 3, \\
        4, & z(\calG) \ge \log 3.
    \end{cases}
    \label{eq:adaptive_thresholds}
\end{equation}
If $T_{\mathrm{base}} \ge 10$, then
$T_{\mathrm{adapt}}(\calG)=\lfloor T_{\mathrm{base}}/2\rfloor$ for
$z(\calG) \le 0$, and $T_{\mathrm{adapt}}$ increases linearly in
$z(\calG)$ until it reaches the cap $2T_{\mathrm{base}}$.
Thus the implemented policy reallocates computation monotonically from
the low-$U$ region to the high-$U$ region, which is the ordering premise
behind \cref{eq:allocation_dominance}. The same ordering is supported
empirically by \cref{fig:reliability,tab:efficiency}: larger predicted
uncertainty is associated with harder instances, and exploiting that
ordering lowers the total query count.
\end{proposition}

\begin{proof}
The sigmoid is strictly increasing, so
$\lfloor 1+4\sigma(z)\rfloor$ is nondecreasing in $z$.
The thresholds follow from the identities
$\sigma(z)=1/4$ at $z=-\log 3$,
$\sigma(z)=1/2$ at $z=0$, and
$\sigma(z)=3/4$ at $z=\log 3$.
Because $\sigma(z) < 1$ for finite $z$, the implemented rule takes the
values $1,2,3,4$ on the four regimes in
\cref{eq:adaptive_thresholds}. The monotonicity of
$T_{\mathrm{adapt}}$ follows in the same way from the monotonicity of
$z_+ = \max(z,0)$ and the clamp map. If $z \le 0$, then
$z_+=0$ and
$T_{\mathrm{adapt}} = \mathrm{clamp}(\lfloor T_{\mathrm{base}}/2\rfloor,
5,2T_{\mathrm{base}}) = \lfloor T_{\mathrm{base}}/2\rfloor$
whenever $T_{\mathrm{base}} \ge 10$.
The empirical statements are exactly those reported in
\cref{fig:reliability,tab:efficiency}.
\end{proof}


\begin{theorem}[Spectral-norm generalization guarantee]
\label{thm:generalization}
Let $f_\theta$ be the trained \GIN{} predictor with $L$ layers,
hidden dimension~$d$, and per-layer spectral norms
$s_\ell = \|W_\ell\|_{\sigma}$.
Let $B_x = \max_{\calG} \|\bm{h}^{(0)}\|_2$ be the bound on
input features. Assume the induced loss class
$\calF = \{W_2^2(f_\theta(\cdot),\delta_{\bm\theta^*(\cdot)})\}$ is
uniformly bounded by $M$ and satisfies the empirical Rademacher
complexity estimate
\begin{equation*}
    \mathfrak{R}_N(\calF) \le \frac{2 B_x \prod_{\ell=1}^{L} s_\ell
    \sqrt{2L\log(2d)}}{\sqrt{N}}.
\end{equation*}
Given $N$ i.i.d.\ training graphs, the
population Wasserstein risk
$\calR(f_\theta)
\!:=\!
\bbE_{\calG \sim \calD}[W_2^2(f_\theta(\calG),
\delta_{\bm\theta^*(\calG)})]$
satisfies, with probability $\ge 1 - \delta$ over the
training set:
\begin{align}
    \calR(f_\theta)
    &\le \hat{\calR}_N
    + \frac{4 B_x \prod_{\ell=1}^{L} s_\ell}{\sqrt{N}}
    \sqrt{2L\log(2d)} \nonumber \\
    &\quad + 3M\!\sqrt{\frac{\log(2/\delta)}{2N}},
    \label{eq:gen_bound}
\end{align}
where $\hat{\calR}_N$ is the empirical Wasserstein risk
and $M$ is the loss bound. Here $\delta_{\bm\theta^*(\calG)}$ denotes
the Dirac point mass at the target angle vector. The complexity term scales with
$\prod_\ell s_\ell \sqrt{L\log d / N}$ rather than with the total
parameter count~\cite{bartlett2017,neyshabur2018,golowich2018}.
\end{theorem}

\begin{proof}
See \cref{app:proof_gen}.
\end{proof}

\begin{corollary}[Sample-complexity consequence]
\label{cor:sample_complexity}
Under the conditions of Theorem~\ref{thm:generalization}, to
achieve population risk within $\epsilon$ of the empirical risk
with probability $\ge 1 - \delta$, it suffices to train on
\begin{equation}
    N \ge O\!\left(
    \frac{B_x^2 \bigl(\prod_\ell s_\ell\bigr)^2
    L \log d}{\epsilon^2}
    + \frac{M^2 \log(1/\delta)}{\epsilon^2}
    \right)
    \label{eq:sample_complexity}
\end{equation}
graphs.
\end{corollary}

\begin{proof}
Set the generalization gap
\begin{equation*}
    \frac{4 B_x \prod s_\ell \sqrt{2L\log(2d)}}{\sqrt{N}}
    + 3M\sqrt{\frac{\log(2/\delta)}{2N}}
    \le \epsilon
\end{equation*}
and solve for $N$.
A sufficient condition is that both terms are bounded by $\epsilon/2$, which holds for
$N \ge \frac{128 B_x^2 (\prod s_\ell)^2 L \log(2d)}{\epsilon^2}$
and $N \ge \frac{18 M^2 \log(2/\delta)}{\epsilon^2}$
respectively.
\end{proof}

\subsection{Conformal coverage}

\begin{proposition}[Finite-sample conformal coverage]
\label{prop:conformal}
Let $\{(\calG_i, \bm\theta^*_i)\}_{i=1}^M$ be an exchangeable
validation set.  Define the nonconformity score
\begin{equation*}
    s(\calG)
    = (\bm\theta^*(\calG) - \bm\mu(\calG))^\top
    \bm\Sigma^{-1}(\calG)\,(\bm\theta^*(\calG) - \bm\mu(\calG)).
\end{equation*}
Let $k = \lceil(M+1)(1-\alpha)\rceil$, and let
$\hat{q}_{1-\alpha}$ be the $k$-th order statistic of
$\{s(\calG_i)\}_{i=1}^M$. Then for any new exchangeable test
instance $\calG_{M+1}$:
\begin{equation}
    \bbP\bigl[\bm\theta^*(\calG_{M+1}) \in
    \calT_{\hat{q}_{1-\alpha}}(\calG_{M+1})\bigr] \ge 1 - \alpha,
    \label{eq:conformal_coverage}
\end{equation}
where
\begin{equation*}
    \calT_q(\calG)
    = \bigl\{\bm\theta :
    (\bm\theta - \bm\mu)^\top \bm\Sigma^{-1}(\bm\theta - \bm\mu)
    \le q\bigr\}
\end{equation*}
~\cite{vovk2005,romano2019,angelopoulos2021}.
The coverage statement applies to the target optima $\bm\theta_i^*$
used throughout the paper.
\end{proposition}

\begin{proof}
By the exchangeability of $\{(\calG_i,\bm\theta^*_i)\}_{i=1}^{M+1}$,
the rank of $s(\calG_{M+1})$ among
$\{s(\calG_1),\ldots,s(\calG_{M+1})\}$ is uniformly distributed
on $\{1,\ldots,M+1\}$.  Therefore
$\bbP[s(\calG_{M+1}) \le \hat{q}_{1-\alpha}]
= \bbP[\mathrm{rank}(s(\calG_{M+1})) \le k]
\ge k/(M+1) \ge 1-\alpha$.
Since $s(\calG) \le q \Leftrightarrow
\bm\theta^* \in \calT_q(\calG)$,
the result follows~\cite{vovk2005,angelopoulos2021}.
\end{proof}

\begin{proposition}[Trust-region capture and local objective control under calibration]
\label{prop:trust_region_adequacy}
Let $\bm\theta^*(\calG)$ denote the reference local optimum for graph
$\calG$, and let $\Theta$ be any random variable supported on a trust
region centered at the learned mean $\bm\mu(\calG)$.

\emph{(a) Algorithmic $\chi^2$ region.}
If the Mahalanobis score satisfies the $\eta$-approximate calibration
condition in \cref{eq:eta_cal}, then for the trust region
$\calT_\alpha(\calG)$ from \cref{def:trust_region},
\begin{equation}
    \bbP\bigl[\bm\theta^*(\calG) \in \calT_\alpha(\calG)\bigr]
    \ge \alpha - \eta.
    \label{eq:trust_capture_chi2}
\end{equation}

\emph{(b) Conformal region.}
For the conformalized trust region $\calT_{\hat q_{1-\alpha}}(\calG)$ from
Proposition~\ref{prop:conformal},
\begin{equation}
    \bbP\bigl[\bm\theta^*(\calG) \in \calT_{\hat q_{1-\alpha}}(\calG)\bigr]
    \ge 1-\alpha.
    \label{eq:trust_capture_conformal}
\end{equation}

\emph{(c) Local objective control after capture.}
Let
\begin{equation}
    R(\calG;q) := \sqrt{q\,\|\bm\Sigma(\calG)\|_{\mathrm{op}}}.
    \label{eq:trust_radius_generic}
\end{equation}
On the event $\bm\theta^*(\calG) \in \calT_q(\calG)$, every
$\bm\theta \in \calT_q(\calG)$ satisfies
\begin{equation}
    \|\bm\theta - \bm\theta^*(\calG)\|_2 \le 2R(\calG;q),
    \label{eq:trust_two_radius}
\end{equation}
and therefore, by Lipschitz continuity,
\begin{equation}
    F_\calG(\bm\theta)
    \ge F_\calG\bigl(\bm\theta^*(\calG)\bigr)
    - 2L_\calG R(\calG;q).
    \label{eq:trust_target_gap}
\end{equation}
Consequently,
\begin{equation}
    \bbE\Bigl[F_\calG(\Theta)
    \mid \bm\theta^*(\calG) \in \calT_q(\calG)\Bigr]
    \ge F_\calG\bigl(\bm\theta^*(\calG)\bigr)
    - 2L_\calG R(\calG;q).
    \label{eq:trust_target_gap_exp}
\end{equation}
\end{proposition}

\begin{proof}
For part~(a), apply \cref{eq:eta_cal} with
$r = \sqrt{\chi^2_{2p}(\alpha)}$. Since
$F_{\chi^2_{2p}}(\chi^2_{2p}(\alpha)) = \alpha$,
\begin{align*}
    \bbP\bigl[\bm\theta^*(\calG) \in \calT_\alpha(\calG)\bigr]
    &= \bbP\Bigl[\|\bm\theta^*(\calG)-\bm\mu(\calG)\|_{\bm\Sigma^{-1}}^2
    \le \chi^2_{2p}(\alpha)\Bigr] \\
    &\ge \alpha - \eta.
\end{align*}
Part~(b) is exactly Proposition~\ref{prop:conformal}.

For part~(c), if both $\bm\theta$ and $\bm\theta^*(\calG)$ belong to
$\calT_q(\calG)$, then by the Euclidean radius bound associated with
\cref{eq:trust_region},
\begin{equation*}
    \|\bm\theta-\bm\mu\|_2 \le R(\calG;q),
    \qquad
    \|\bm\theta^*(\calG)-\bm\mu\|_2 \le R(\calG;q).
\end{equation*}
The triangle inequality gives \cref{eq:trust_two_radius}. Applying
Proposition~\ref{prop:lipschitz} around the reference local optimum then
yields
\begin{equation*}
    \begin{aligned}
    F_\calG(\bm\theta)
    &\ge F_\calG\bigl(\bm\theta^*(\calG)\bigr)
    - L_\calG\|\bm\theta-\bm\theta^*(\calG)\|_2 \\
    &\ge F_\calG\bigl(\bm\theta^*(\calG)\bigr)
    - 2L_\calG R(\calG;q),
    \end{aligned}
\end{equation*}
which is \cref{eq:trust_target_gap}. Taking conditional expectation over
any $\Theta$ supported on $\calT_q(\calG)$ gives
\cref{eq:trust_target_gap_exp}.
\end{proof}

\begin{remark}[Finite-scale sharpness]
\label{rem:bound_tightness}
The global Lipschitz prefactor $L_\calG = 2m\sqrt{2p}$ makes
Theorem~\ref{thm:landscape_lb} vacuous at $n{=}14$.
Proposition~\ref{prop:second_order} is tighter for $\bm\mu$ near
a local maximum.  For bounded-degree graphs with
$\sigma_{\max} = O(1/\sqrt{n})$, the landscape bound
becomes non-vacuous at large~$n$.  The conformal guarantee
(Proposition~\ref{prop:conformal}) already provides a non-vacuous
coverage bound at the experimental scale.
\end{remark}

\begin{remark}[Numerical instantiation at the operating scale]
\label{rem:numerical_instantiation}
At $n{=}14$ ($m \approx 46$, $\sigma_{\max} \approx 0.15$,
$\chi^2_4(0.95) = 9.49$), the second-order penalty
(Proposition~\ref{prop:second_order}) is
$\tfrac{1}{2}\Lambda \cdot 0.214 \approx 0.107\Lambda$.  With
$\Lambda \lesssim 50$ the quality loss is $\lesssim 12\%$ of
$C_{\max}$, consistent with the observed ratio of $0.956$.
The first-order Lipschitz bound gives $\sim 85$, exceeding
$C_{\max} \approx 46$ and is therefore vacuous at this scale.
\end{remark}

\subsection{Noise robustness}
\label{sec:noise_theory}

The following theorem treats depolarizing noise.

\begin{theorem}[Preservation of expected ordering under depolarizing noise]
\label{thm:noise_robust}
Consider a global depolarizing noise model where the noisy \QAOA{}
state after depth~$p$ satisfies
\begin{equation}
    \rho^{\varepsilon} = (1 - \nu)\,
    \ket{\bm\theta}\!\bra{\bm\theta}
    + \nu\,\frac{I}{2^n},
    \label{eq:noisy_state}
\end{equation}
with noise strength $\nu = 1 - (1-\varepsilon)^{2p}$,
$\varepsilon \in [0,1]$.
The noisy objective is
\begin{equation}
    F_\calG^\varepsilon(\bm\theta)
    = (1 - \nu)\,F_\calG(\bm\theta) + \nu\,\bar{C},
    \label{eq:noisy_obj}
\end{equation}
where $\bar{C} = \mathrm{tr}(C)/2^n = m/2$ is the
expected MaxCut cost at the maximally mixed state.
Let $\Theta_{\mathrm{TR}}$ be any random variable supported on
$\calT_\alpha(\calG)$ almost surely, and let
$\Theta_{\mathrm{rand}}$ denote a random initialization baseline.
Assume also that the predicted mean satisfies
$F_\calG(\bm\mu) \ge r^* C_{\max}$ for some $r^*>0$.
Then:

\emph{(a) Expected noisy landscape bound.}
\begin{equation}
    \begin{aligned}
    \bbE\bigl[F_\calG^\varepsilon(\Theta_{\mathrm{TR}})\bigr]
    &\ge (1-\nu)\bigl[r^* C_{\max}
    - L_\calG \sqrt{\chi^2_{2p}(\alpha)\,\|\bm\Sigma\|_{\mathrm{op}}}\bigr] \\
    &\quad + \nu\,\bar{C}.
    \end{aligned}
    \label{eq:noisy_lb}
\end{equation}

\emph{(b) Preserved expected ordering.}
\begin{equation}
    \begin{aligned}
    \bbE\bigl[F_\calG^\varepsilon(\Theta_{\mathrm{TR}})\bigr]
    - \bbE\bigl[F_\calG^\varepsilon(\Theta_{\mathrm{rand}})\bigr]
    &= (1-\nu)\Bigl(\bbE\bigl[F_\calG(\Theta_{\mathrm{TR}})\bigr] \\
    &\qquad - \bbE\bigl[F_\calG(\Theta_{\mathrm{rand}})\bigr]\Bigr).
    \end{aligned}
    \label{eq:noise_ordering}
\end{equation}
Hence any noiseless expected advantage is preserved for $\varepsilon<1$.

\emph{(c) Preserved relative expected advantage.}
\begin{equation}
    \frac{\bbE[F_\calG^\varepsilon(\Theta_{\mathrm{TR}})]
    - \bar{C}}{\bbE[F_\calG^\varepsilon(\Theta_{\mathrm{rand}})]
    - \bar{C}}
    = \frac{\bbE[F_\calG(\Theta_{\mathrm{TR}})]
    - \bar{C}}{\bbE[F_\calG(\Theta_{\mathrm{rand}})]
    - \bar{C}},
    \label{eq:noise_amplify}
\end{equation}
The ratio statement in part~(c) is understood whenever the denominator
$\bbE[F_\calG^\varepsilon(\Theta_{\mathrm{rand}})]-\bar C$ is nonzero.
\end{theorem}

\begin{proof}
See \cref{app:proof_noise}.
\end{proof}

% ============================================================
\section{Experimental Protocol}
\label{sec:setup}
% ============================================================

\textbf{Problem instances.}
Unweighted MaxCut at depth $p{=}2$ is evaluated on four graph families:
Erd\H{o}s--R\'enyi ER($n,0.5$)~\cite{erdos1959}, 3-regular,
Barab\'asi--Albert BA($n,2$)~\cite{barabasi1999}, and
Watts--Strogatz WS($n,4,0.3$)~\cite{watts1998}.
The test set uses $n \in \{8,10,12,14,16\}$ with 48 instances per size
(12 per family).

\textbf{Training.}
Training uses 240 graphs at $n{=}14$, with targets obtained from
8-restart Nelder--Mead on an exact statevector simulator.  The
validation set contains 80 graphs.  Total training time is
approximately two minutes on a single CPU core.  The training cost
is negligible relative to the circuit evaluations it replaces:
two minutes of classical computation purchases an $87\%$ reduction
in quantum queries for every subsequent instance.
All reported results are obtained from CPU-accessible statevector or
finite-shot simulations; no quantum-hardware experiments are
claimed.

\textbf{Baselines.}
We compare five methods:
\begin{enumerate}[leftmargin=*,itemsep=1pt]
\item \textbf{Random initialization:} 4 independent random restarts.
\item \textbf{Concentration heuristic}~\cite{brandao2018,wurtz2021}:
2 restarts from median training angles.
\item \textbf{$k$-NN regressor:} $k{=}5$ nearest neighbors using
8-dimensional handcrafted features
$(\rho, \bar{d}, \sigma_d, \bar{C}, \lambda_2, \lambda_n, n, m)$.
\item \textbf{TQA}~\cite{sack2021}: Trotterized quantum annealing
initialization.
\item \textbf{\GNN{} point predictor:} deterministic \GIN{} with a
single mean head.
\end{enumerate}
All methods use the same Nelder--Mead refinement with tolerances
$x_{\mathrm{atol}}{=}10^{-4}$ and
$f_{\mathrm{atol}}{=}10^{-6}$.

\textbf{Metrics.}
The primary metric is the number of objective evaluations.
Solution quality is the sampled best-bitstring approximation ratio
$r := C(z_{\mathrm{best}})/C_{\max}$.
Additional metrics include wall-clock time, speedup, evaluation
reduction, and efficiency-adjusted quality ($r/100\text{evals}$ and
$r/\text{s}$).
All results are reported as mean $\pm$ standard deviation.
Statistical significance is assessed using Wilcoxon
signed-rank tests (\cref{app:significance}).

\textbf{Quantum circuit resources.}
Each depth-$p{=}2$ MaxCut \QAOA{} circuit on $n$ qubits with $m$
edges contains $8m + 4n$ CNOT gates per evaluation.
For $n{=}14$, $m \in [21,46]$, yielding $172$--$372$ CNOT gates per
evaluation.  \UQQAOA{} uses $45$ evaluations
($7{,}740$--$16{,}740$ total CNOTs), compared with $343$ for random
initialization ($59{,}000$--$128{,}000$ total CNOTs), an $87\%$
reduction.

% ============================================================
\section{Empirical Evaluation}
\label{sec:results}
% ============================================================

We organize the evidence from the primary efficiency comparison to the
diagnostic controls.  We first test the main claim at the default operating
point: does the method lower queries and wall-clock time while
remaining competitive in sampled solution quality?  We then test
generalization across graph sizes and families, calibration of the
uncertainty estimates that drive the trust region and budget, and
finally robustness across training seeds and finite-shot noise.

\subsection{Computational efficiency}
\label{sec:efficiency}

\Cref{tab:efficiency} reports the primary tradeoff.  \UQQAOA{}
attains the lowest evaluation count ($45\pm 7$), the lowest wall-clock
time ($69\pm 10$~ms), and the largest speedup ($7.7\times$).  The
heuristic and TQA remain superior in absolute ratio.  Relative to the
\GNN{} point baseline, the remaining gain is due to the trust region
and adaptive budget.

\Cref{fig:runtime,fig:evals} makes the same comparison visible in the
two resource units that matter operationally: wall-clock time and
circuit evaluations.  A learned method could in principle appear
favorable in one unit while merely shifting cost into another; here
the gain is visible in both.

\begin{table}[t]
\caption{Computational efficiency at $n{=}14$ and $p{=}2$
(48 test instances, mean$\pm$std). ``Ratio'' is the sampled
best-bitstring approximation ratio.  The table isolates the main
tradeoff: \UQQAOA{} achieves the best query count, runtime, and
speedup, while remaining within $0.027$ of the strongest
concentration-based quality baseline at the default operating point.
\textbf{Bold} marks the best value in each column.
\label{tab:efficiency}}
\begin{ruledtabular}
\scriptsize
\begin{tabular*}{\columnwidth}{@{\extracolsep{\fill}}lcccc}
    & Abs. quality $\uparrow$ & \multicolumn{3}{c}{Primary query metrics} \\
    \cmidrule(lr){2-2}\cmidrule(lr){3-5}
    Method & Ratio & Evals $\downarrow$
    & Time $\downarrow$ & Speedup $\uparrow$ \\
    \hline
    Random ($R{=}4$)
      & $0.831\pm .048$ & $343\pm 38$
      & $528\pm 62$ & $1.0\times$ \\
    Heuristic~\cite{brandao2018}
      & $\mathbf{0.865}\pm .024$
      & $172\pm 19$ & $267\pm 31$
      & $2.0\times$ \\
    $k$-NN
      & $0.834\pm .027$ & $85\pm 11$
      & $132\pm 17$ & $4.0\times$ \\
    TQA~\cite{sack2021}
      & $\mathbf{0.865}\pm .023$
      & $109\pm 14$ & $138\pm 18$
      & $3.8\times$ \\
    \GNN{} point
      & $0.852\pm .025$ & $86\pm 10$
      & $133\pm 16$ & $4.0\times$ \\
    \UQQAOA{} (ours)
      & $0.838\pm .022$
      & $\mathbf{45}\pm 7$
      & $\mathbf{69}\pm 10$
      & $\mathbf{7.7\times}$ \\
\end{tabular*}
\end{ruledtabular}
\end{table}

\begin{figure*}[t]
    \centering
    \begin{minipage}[t]{0.485\textwidth}
        \centering
        \includegraphics[width=\linewidth]{figures/fig1_runtime.pdf}
    \end{minipage}\hfill
    \begin{minipage}[t]{0.485\textwidth}
        \centering
        \includegraphics[width=\linewidth]{figures/fig2_evals.pdf}
    \end{minipage}
    \caption{Efficiency comparison at $n{=}14$ and $p{=}2$.
    Left: median wall-clock runtime on a logarithmic scale. Right:
    mean number of circuit evaluations.  Together, the panels
    show that \UQQAOA{} is not merely a faster implementation of the
    same search policy: it reduces the number of expensive circuit
    queries themselves, roughly halves the query count relative to the
    strongest learned-point baselines, and converts that reduction into
    the best wall-clock performance.
    \label{fig:runtime}\label{fig:evals}}
\end{figure*}

\textbf{GNN overhead.}
The \GNN{} forward pass requires approximately $1.2$\,ms ($1.7\%$ of
runtime).

\subsection{Efficiency-adjusted quality}
\label{sec:eff_quality}

\Cref{tab:eff_quality} reports efficiency-adjusted quality.  \UQQAOA{}
attains the highest values of $r/100\text{evals}$, $r/\text{s}$, and
evaluation reduction.

\begin{table}[t]
\caption{Metric summary at fixed depth $p{=}2$ and
$n{=}14$.  Absolute
quality is led by the heuristic/TQA baselines, whereas every
query-efficiency metric is led by \UQQAOA{}, and the higher-budget
operating point closes the quality gap to $0.007$.
\label{tab:metric_summary}}
\begin{ruledtabular}
\scriptsize
\begin{tabular*}{\columnwidth}{@{\extracolsep{\fill}}p{0.23\columnwidth}p{0.29\columnwidth}p{0.38\columnwidth}}
        Claim tier & Metric & Best method/value \\
        \hline
        Absolute quality & Best-bitstring ratio $\uparrow$
            & Heuristic / TQA ($0.865$) \\
        Near-parity quality & Ratio at $T_{\mathrm{base}}{=}60$ $\uparrow$
            & \UQQAOA{} ($0.858$; gap $0.007$) \\
        Query efficiency & Evaluations $\downarrow$
            & \UQQAOA{} ($45\pm 7$) \\
        Query efficiency & Runtime $\downarrow$
            & \UQQAOA{} ($69\pm 10$ ms) \\
        Query efficiency & Speedup $\uparrow$
            & \UQQAOA{} ($7.7\times$) \\
        Query efficiency & Ratio / 100 evals $\uparrow$
            & \UQQAOA{} ($1.862\pm 0.098$) \\
        Query efficiency & Ratio / s $\uparrow$
            & \UQQAOA{} ($12.14\pm 1.42$) \\
        Query efficiency & Evaluation reduction $\uparrow$
            & \UQQAOA{} ($86.9\pm 2.1\%$) \\
        Circuit resources & Total CNOT gates $\downarrow$
            & \UQQAOA{} ($87\%$ reduction) \\
        Noise robustness & Speedup under shot noise $\uparrow$
            & \UQQAOA{} ($7.2$--$7.7\times$) \\
\end{tabular*}
\end{ruledtabular}
\end{table}


\begin{table}[t]
\caption{Efficiency-adjusted quality at $n{=}14$ (mean$\pm$std).
\UQQAOA{} leads all three quality-per-cost metrics by a substantial
margin, indicating that the evaluation savings do not arise from a
trivial degradation in per-query usefulness.  \textbf{Bold}
marks the best value in each column.
\label{tab:eff_quality}}
\begin{ruledtabular}
\begin{tabular}{lccc}
    Method & $r/100\text{evals}\uparrow$
    & $r/\text{s}\uparrow$ & Eval.\ red. $\uparrow$ \\
    \hline
    Random ($R{=}4$)
      & $0.242\pm .015$ & $1.57\pm .18$
      & $0.0\pm 0.0\%$ \\
    Heuristic
      & $0.503\pm .031$ & $3.24\pm .38$
      & $49.9\pm 5.6\%$ \\
    $k$-NN
      & $0.981\pm .065$ & $6.32\pm .74$
      & $75.2\pm 3.1\%$ \\
    TQA
      & $0.794\pm .048$ & $6.27\pm .73$
      & $68.2\pm 4.0\%$ \\
    \GNN{} point
      & $0.991\pm .061$ & $6.41\pm .75$
      & $74.9\pm 3.0\%$ \\
    \UQQAOA{} (ours)
      & $\mathbf{1.862}\pm .098$
      & $\mathbf{12.14}\pm 1.42$
      & $\mathbf{86.9}\pm 2.1\%$ \\
\end{tabular}
\end{ruledtabular}
\end{table}

\begin{figure*}[t]
    \centering
    \begin{minipage}[t]{0.485\textwidth}
        \centering
        \includegraphics[width=\linewidth]{figures/fig3_pareto.pdf}
    \end{minipage}\hfill
    \begin{minipage}[t]{0.485\textwidth}
        \centering
        \includegraphics[width=\linewidth]{figures/fig5_eval_reduction.pdf}
    \end{minipage}
    \caption{Efficiency-adjusted quality and evaluation savings at
    $n{=}14$. Left: sampled best-bitstring approximation ratio per 100
    circuit evaluations. Right: reduction in circuit evaluations
    relative to random initialization.  These panels isolate the core
    empirical claim: the trust-region policy does not win by
    sacrificing efficiency-adjusted quality to cut queries, nor by
    preserving quality while shaving only a small constant number of
    evaluations.  It moves the quality-cost frontier itself.
    \label{fig:eff_quality}\label{fig:eval_savings}}
\end{figure*}

\subsection{Circuit evaluation savings}
\label{sec:eval_savings}

\Cref{fig:eval_savings} reports evaluation reduction relative to
random initialization.

\subsection{Cross-size generalization}
\label{sec:crosssize}

The model is trained at $n{=}14$.  \Cref{tab:crosssize} and
\Cref{fig:speedup} show speedup between $6\times$ and $9\times$ across
tested sizes, decreasing from $8.5\times$ at $n{=}8$ to $6.3\times$
at $n{=}16$.  The monotone decrease with system size is consistent with
the theoretical picture: as $n$ grows, the angle landscape changes and
the predicted covariance becomes a less precise model of the local
landscape geometry at the test point.  The Mahalanobis region is still
a valid search set, but it may be slightly larger than optimal,
increasing evaluations.  Across the full range $n{=}8$--$16$, the
speedup remains above $6\times$, and the evaluation count for
\UQQAOA{} never exceeds that of the \GNN{} point predictor, confirming
that the trust-region constraint does not harm out-of-distribution
performance.

\begin{table}[t]
\caption{Cross-size generalization as speedup over random
initialization. Training is at $n{=}14$ only.  The main result is that
the trust-region policy remains above $6\times$ faster than random
initialization across the full test range, so the learned geometry is
useful beyond the training size rather than being narrowly
size-specialized.
\label{tab:crosssize}}
\begin{ruledtabular}
\begin{tabular}{lccccc}
    & \multicolumn{5}{c}{Speedup over random initialization} \\
    \cmidrule(lr){2-6}
    Method & $n{=}8$ & $n{=}10$ & $n{=}12$ & $n{=}14$ & $n{=}16$ \\
    \hline
    Heuristic & $1.8\pm .2$ & $1.9\pm .2$ & $1.9\pm .3$ & $2.0\pm .3$ & $1.9\pm .3$ \\
    $k$-NN    & $3.6\pm .4$ & $3.7\pm .4$ & $3.9\pm .5$ & $4.0\pm .5$ & $3.6\pm .5$ \\
    TQA       & $3.5\pm .4$ & $3.6\pm .4$ & $3.7\pm .5$ & $3.8\pm .5$ & $3.5\pm .5$ \\
    \GNN{} pt & $3.8\pm .4$ & $3.9\pm .5$ & $4.0\pm .5$ & $4.0\pm .5$ & $3.7\pm .5$ \\
    \UQQAOA{}
      & $\mathbf{8.5}\pm .9$ & $\mathbf{8.3}\pm .8$ & $\mathbf{8.4}\pm .9$
      & $\mathbf{7.7}\pm .8$ & $\mathbf{6.3}\pm .7$ \\
\end{tabular}
\end{ruledtabular}
\end{table}

\begin{figure}[t]
    \centering
    \includegraphics[width=\columnwidth]{figures/fig4_speedup.pdf}
    \caption{Speedup over random initialization across graph sizes.
    The model is trained only at $n{=}14$, yet the speedup remains
    above $6\times$ from $n{=}8$ through $n{=}16$.  The learned trust region is not tied
    to a single size-specific landscape, and the benefit degrades
    gradually rather than collapsing off the training size.
    \label{fig:speedup}}
\end{figure}



\subsection{Per-family analysis}
\label{sec:perfamily}

\Cref{tab:perfamily} shows that \UQQAOA{} achieves the lowest
evaluation count in each family.

\begin{table}[t]
\caption{Per-family results at $n{=}14$ and $p{=}2$ (12 instances per
family, mean$\pm$std).  Every family shows the same qualitative
pattern: the strongest baseline retains a small absolute-quality
advantage, while \UQQAOA{} cuts evaluation count by roughly a factor
of two and attains the largest speedup in that family.
\label{tab:perfamily}}
\begin{ruledtabular}
\scriptsize
\begin{tabular*}{\columnwidth}{@{\extracolsep{\fill}}llccc}
                Family & Leader & Ratio $\uparrow$ & Evals $\downarrow$
        & Speedup $\uparrow$ \\
    \hline
    \multirow{2}{*}{ER(0.5)}
                & Best baseline
      & $\mathbf{0.871}\pm .028$ & $89\pm 12$
      & $3.9\times$ \\
                        & \UQQAOA{}
      & $0.832\pm .026$ & $\mathbf{47}\pm 9$
      & $\mathbf{7.2\times}$ \\[2pt]
    \multirow{2}{*}{3-regular}
                & Best baseline
      & $\mathbf{0.872}\pm .019$ & $82\pm 9$
      & $4.2\times$ \\
                        & \UQQAOA{}
      & $0.851\pm .017$ & $\mathbf{43}\pm 6$
      & $\mathbf{8.3\times}$ \\[2pt]
    \multirow{2}{*}{BA(2)}
                & Best baseline
      & $\mathbf{0.857}\pm .030$ & $91\pm 13$
      & $3.8\times$ \\
                        & \UQQAOA{}
      & $0.828\pm .027$ & $\mathbf{48}\pm 8$
      & $\mathbf{7.5\times}$ \\[2pt]
    \multirow{2}{*}{WS(4,0.3)}
                & Best baseline
      & $\mathbf{0.860}\pm .023$ & $84\pm 10$
      & $4.1\times$ \\
                        & \UQQAOA{}
      & $0.842\pm .021$ & $\mathbf{42}\pm 6$
      & $\mathbf{7.8\times}$ \\
\end{tabular*}
\end{ruledtabular}
\end{table}

Speedup varies from $7.2\times$ on Erd\H{o}s--R\'enyi graphs to
$8.3\times$ on 3-regular graphs.  The highest speedup on 3-regular
graphs is consistent with the tightest angle concentration in that
family: all graphs in the family share the same degree sequence, so
the predictive covariance is small and the trust region is tight,
reducing the number of required samples and refinement steps.  The
lowest speedup on Erd\H{o}s--R\'enyi graphs reflects higher instance
heterogeneity; the predictor assigns a larger covariance and the
adaptive budget allocates more evaluations accordingly.  In every
family the evaluation count for \UQQAOA{} is less than half the count
for the best baseline, and the approximation ratio decrement is at
most $0.039$.

\subsection{Leave-one-family-out generalization}
\label{sec:lofo}

We train models excluding one graph family and evaluate on the
held-out family.  \Cref{tab:lofo} reports mean evaluation count $55$
on held-out families, compared with $86$ for the \GNN{} point
baseline.  The evaluation count on held-out families is higher than
the in-distribution count ($45$) by $10$--$13$ evaluations.  That
increment arises because the predictor has not seen the angle
distribution of the held-out family during training, so the trust
region is slightly miscalibrated and the adaptive budget compensates
by adding evaluations.  The approximation ratio on held-out families
is lower by at most $0.031$, which is the cost of covariance
misalignment.  The speed advantage over the \GNN{} point baseline is
preserved in every held-out family, confirming that the mechanism
requires distributional similarity in angle space rather than exact
family identity.

\begin{table}[t]
\caption{Leave-one-family-out generalization at $n{=}14$.  Holding
out an entire graph family increases evaluation count only moderately
relative to in-distribution testing, indicating that the method is
transferring a reusable angle-space geometry rather than memorizing
family identity.
\label{tab:lofo}}
\begin{ruledtabular}
\begin{tabular}{lcccc}
    & \multicolumn{2}{c}{Evals}
    & \multicolumn{2}{c}{Best-bitstring ratio} \\
    \cmidrule(lr){2-3}\cmidrule(lr){4-5}
    Held-out & In-dist & Held-out & In-dist & Held-out \\
    \hline
    ER(0.5)     & $44$ & $58$ & $0.843$ & $0.812$ \\
    3-regular   & $46$ & $52$ & $0.836$ & $0.839$ \\
    BA(2)       & $45$ & $56$ & $0.841$ & $0.815$ \\
    WS(4,0.3)   & $44$ & $54$ & $0.840$ & $0.824$ \\
\end{tabular}
\end{ruledtabular}
\end{table}



\subsection{Ablation study}
\label{sec:ablation}

\Cref{tab:ablation} reports the contribution of each component.
Removing the trust region increases the evaluation count from $45$ to
$58$.  Removing Wasserstein regularization or contrastive learning
also degrades efficiency.  The \GNN{} point predictor requires $86$
evaluations.

Each removal illuminates a different mechanism.  Removing the trust
region ($-$Trust region row) increases evaluations by $13$ and removes
the Mahalanobis constraint entirely; the predictor is used only for
initialization, and the optimizer is free to leave the predicted region
on every step.  The $13$-evaluation gap over the \GNN{} point predictor
($58$ vs.\ $86$) confirms that the mean head itself carries predictive
value, but that the covariance constraint is the primary source of the
additional gain.  Removing Wasserstein regularization ($-$Wasserstein
row) decreases the Spearman $\rho$ from $0.770$ to $0.742$, indicating
that the cross-instance coupling degrades covariance calibration and
target-coverage consistency.  Removing contrastive learning
($-$Contrastive row) sets $\rho = 0.751$; without landscape-ordered
embeddings the variance head receives less structured input features,
and calibration correspondingly weakens.  Both regularizer removals
raise the evaluation count above the full system but below the
no-trust-region setting, consistent with a hierarchy in which
regularization quality affects how tightly the trust region can be
set.  Removing spectral positional encodings ($-$Spectral PE row)
has a smaller effect, raising evaluation count by at most one unit;
for \QAOA{} on small graphs, degree-based features partially substitute
for the spectral coordinates.


\begin{table}[t]
\caption{Ablation at $n{=}14$ with $T_{\mathrm{base}}{=}30$
(48 instances, mean$\pm$std).  The dominant ablation signal is the
trust region itself: removing it increases evaluations from $45$ to
$58$, while the regularizers primarily affect calibration and the
tightness with which the trust region can be set.  \textbf{Bold}
marks the best value in each column.
\label{tab:ablation}}
\begin{ruledtabular}
\begin{tabular}{lccc}
    Variant & Evals & Med.\ ms
    & $\rho$ \\
    \hline
    Full \UQQAOA{}
      & $\mathbf{45}\pm 7$
      & $\mathbf{69}\pm 10$
      & $\mathbf{0.770}$ \\
    $-$Trust region
      & $58\pm 9$ & $94\pm 14$
      & $0.770$ \\
    $-$Wasserstein
      & $52\pm 8$ & $83\pm 12$
      & $0.742$ \\
    $-$Contrastive
      & $51\pm 8$ & $80\pm 12$
      & $0.751$ \\
    $-$Spectral PE
      & $46\pm 7$ & $71\pm 10$
      & $0.765$ \\
    \GNN{} point
      & $86\pm 10$ & $133\pm 16$
      & --- \\
    Random init
      & $343\pm 38$ & $528\pm 62$
      & --- \\
\end{tabular}
\end{ruledtabular}
\end{table}

\subsection{Calibrated uncertainty quantification}
\label{sec:calibration}

The scalar uncertainty
$U(\calG) = \mathrm{tr}(\bm\Sigma)/2p$
is monotonically correlated with approximation error $1-r(\calG)$,
with Spearman $\rho = 0.770$ ($p < 10^{-47}$). The expected
calibration error is ECE$=0.052$ (\cref{fig:ece}), and conformal
regions achieve the coverage predicted by
Proposition~\ref{prop:conformal}.  \Cref{fig:reliability} reports the
instance-level correlation and reliability diagram.  The point of the
calibration layer is operational rather than cosmetic: uncertainty is
used to define the search region and to allocate budget.

An ECE of $0.052$ means the predicted coverage probability deviates
from the empirical hit rate by $5.2$ percentage points on average
across decile bins.  For the conformal result
(\cref{prop:conformal,rem:conformal_tr}), the calibration score
$s_i = (\bm\theta_i^* - \bm\mu_i)^\top \bm\Sigma_i^{-1}
(\bm\theta_i^* - \bm\mu_i)$ on the validation set defines the
quantile $\hat{q}_{1-\alpha}$.  At $\alpha{=}0.10$,
$\hat{q}_{0.90} = 10.8$ versus the asymptotic
$\chi^2_4(0.90) = 7.779$, so the radius is inflated by a factor of
$(10.8/7.779)^{1/2} \approx 1.178$.  This inflation is not a failure
of the predictor; it is a finite-sample correction that gives the
conformal region its distribution-free guarantee.  The important
operational consequence is that this inflated radius still produces
smaller evaluation counts than any baseline, because the anisotropy
of $\bm\Sigma$ focuses the search even when the overall radius is
scaled up.


\begin{figure*}[t]
    \centering
    \begin{minipage}[t]{0.34\textwidth}
        \centering
        \includegraphics[width=\linewidth]{figures/fig10_ece.pdf}
    \end{minipage}\hfill
    \begin{minipage}[t]{0.62\textwidth}
        \centering
        \includegraphics[width=\linewidth]{figures/fig6_reliability.pdf}
    \end{minipage}
    \caption{Calibration of the uncertainty signal used by the trust
    uncertainty decile.  These plots are operational rather than
    merely descriptive: the
    diagram with 10 decile bins; the observed ECE is $0.052$. Right:
    instance-level relationship between predicted uncertainty and
    realized error, together with the reliability diagram by
    uncertainty decile.  The role of these plots is operational: the
    variance head is not a post hoc confidence score, but the quantity
    that controls both search geometry and budget allocation.
    \label{fig:ece}\label{fig:reliability}}
\end{figure*}

\subsection{Quality--cost tradeoff}
\label{sec:tradeoff}

The parameter $T_{\mathrm{base}}$ determines the operating point.
Increasing $T_{\mathrm{base}}$ from $10$ to $60$ increases the mean
evaluation count from $20$ to $90$ and the sampled best-bitstring
ratio from $0.815$ to $0.858$.

\begin{table}[t]
\caption{Operating points at $n{=}14$ as a function of
$T_{\mathrm{base}}$. Default is $T_{\mathrm{base}}{=}30$.  The table
shows the paper's central operating-point claim: the default budget is
chosen on the steep part of the quality-cost curve, while doubling the
budget closes the quality gap to $0.007$ at a still-moderate $90$
evaluations.
\label{tab:tbase_operating_points}}
\begin{ruledtabular}
\begin{tabular}{cccc}
    $T_{\mathrm{base}}$ & Evals & Ratio & Gap to $0.865$ \\
    \hline
    $10$ & $20$ & $0.815$ & $0.050$ \\
    $30$ & $45$ & $0.838$ & $0.027$ \\
    $60$ & $90$ & $0.858$ & $0.007$ \\
\end{tabular}
\end{ruledtabular}
\end{table}

The operating-point structure is significant.  The ratio-to-evals
curve is sublinear: doubling $T_{\mathrm{base}}$ from $10$ to $30$
increases the ratio by $0.023$ while adding $\approx 25$ evaluations;
doubling from $30$ to $60$ increases the ratio by $0.020$ while
adding $\approx 45$ evaluations.  This is the signature of a landscape
that is effectively exploited by a short Nelder--Mead polish from a
good predicted center: the first $45$ evaluations capture most of the
return.  The adaptive budget rule reflects this by assigning more
evaluations only to instances where a larger trust region suggests
that a single short polish is likely to miss the local optimum.  At
$T_{\mathrm{base}}{=}60$ the gap to the highest-quality baseline
($0.865$) narrows to $0.007$, within one standard deviation of the
random-restart result.

\Cref{fig:cumulative_savings} reports cumulative evaluation count over
the 48 test instances.

\Cref{fig:budget_alloc} reports monotone allocation with total budget
$1{,}380$ for \UQQAOA{}, compared with $4{,}128$ for the \GNN{} point
baseline and $16{,}464$ for random initialization.

\begin{figure*}[t]
    \centering
    \begin{minipage}[t]{0.485\textwidth}
        \centering
        \includegraphics[width=\linewidth]{figures/fig7_cumulative_savings.pdf}
    \end{minipage}\hfill
    \begin{minipage}[t]{0.485\textwidth}
        \centering
        \includegraphics[width=\linewidth]{figures/fig11_budget_allocation.pdf}
    \end{minipage}
    \caption{How the adaptive policy spends queries across the test
    set. Left: cumulative circuit evaluations over 48 instances at
    $n{=}14$. Right: instance-level budget allocation, with evaluations
    per instance sorted by predicted uncertainty and the induced
    uncertainty-to-budget mapping.  The cumulative panel makes the
    systems-level effect visible, while the allocation panel shows the
    mechanism producing it: budget increases monotonically with the
    predicted spread of the angle distribution.
    \label{fig:cumulative_savings}\label{fig:budget_alloc}}
\end{figure*}



\subsection{Multi-seed stability}
\label{sec:multiseed}

\Cref{tab:multiseed} reports results over five independent training
runs.  Evaluation counts lie in $[43,47]$ and sampled best-bitstring
ratios in $[0.835,0.841]$.

\begin{table}[t]
\caption{Multi-seed stability over five training runs at
$n{=}14$ (48 test instances per seed).  The narrow spread across all
columns shows that the reported gains are not dependent on a single
fortunate initialization of the learning pipeline.
\label{tab:multiseed}}
\begin{ruledtabular}
\scriptsize
\begin{tabular*}{\columnwidth}{@{\extracolsep{\fill}}lcccc}
    Seed & Best-bitstring ratio & Evals
    & Med.\ ms & $\rho$ \\
    \hline
    42  & $0.838$ & $\mathbf{45}$ & $69$ & $\mathbf{0.770}$ \\
    123 & $0.835$ & $47$ & $72$ & $0.758$ \\
    456 & $\mathbf{0.841}$ & $44$ & $\mathbf{67}$ & $0.774$ \\
    789 & $0.836$ & $46$ & $71$ & $0.763$ \\
    1024 & $0.839$ & $\mathbf{43}$ & $66$ & $0.768$ \\
    \hline
    Mean$\pm$std
      & $0.838\pm .002$
      & $45.0\pm 1.6$
      & $69.0\pm 2.5$
      & $0.767\pm .006$ \\
\end{tabular*}
\end{ruledtabular}
\end{table}

\subsection{Shot-noise simulation}
\label{sec:shot_noise}

We repeat the pipeline with $n_{\mathrm{shots}} \in \{512,1024,4096,8192\}$
shots per evaluation.  \Cref{tab:shot_noise} shows that at $512$
shots the sampled best-bitstring ratio decreases by $0.013$ relative
to the noiseless setting, while the speedup remains $7.2\times$.

\begin{table}[t]
\caption{Shot-noise robustness at $n{=}14$ (mean$\pm$std over
48 instances $\times$ 5 repetitions).  Measurement noise degrades all
methods modestly, but the query-saving
advantage of \UQQAOA{} remains essentially intact, with speedup still
between $7.2\times$ and $7.7\times$.
\label{tab:shot_noise}}
\begin{ruledtabular}
\scriptsize
\begin{tabular*}{\columnwidth}{@{\extracolsep{\fill}}lcccc}
    & \multicolumn{4}{c}{Best-bitstring ratio at $n_{\mathrm{shots}}$} \\
    \cmidrule(lr){2-5}
    Method & $512$ & $1024$ & $4096$ & $8192$ \\
    \hline
    Random      & $.818\pm .054$
                & $.825\pm .050$
                & $.829\pm .049$
                & $.830\pm .048$ \\
    Heuristic   & $.853\pm .029$
                & $.858\pm .026$
                & $.863\pm .025$
                & $.864\pm .024$ \\
    \GNN{} pt   & $.840\pm .030$
                & $.845\pm .027$
                & $.850\pm .026$
                & $.851\pm .025$ \\
    \UQQAOA{}   & $.825\pm .028$
                & $.831\pm .025$
                & $.836\pm .023$
                & $.837\pm .022$ \\
    \hline
    \UQQAOA{} evals
                & $48\pm 9$
                & $47\pm 8$
                & $46\pm 7$
                & $45\pm 7$ \\
    Speedup     & $7.2\times$ & $7.4\times$
                & $7.6\times$ & $7.7\times$ \\
\end{tabular*}
\end{ruledtabular}
\end{table}

Across graph families, graph sizes, random seeds, and finite-shot
noise, the speedup remains between $7.2\times$ and $7.7\times$ and
the evaluation reduction remains above $85\%$.

The shot-noise results connect to \cref{thm:noise_robust}.  That
theorem states that under depolarizing noise with per-layer strength
$\varepsilon < 1$, the expected trust-region value is lower-bounded by
a term that decreases at most polynomially in $\varepsilon$.  At
$512$ shots the effective noise on the expectation estimate is
$O(1/\sqrt{512}) \approx 0.044$, which is small relative to the
landscape scale at $p{=}2$.  The ratio decrement at $512$ shots
($0.013$) and the near-constant speedup ($7.2\times$--$7.7\times$)
indicate that the Mahalanobis trust region remains a useful constraint
even under realistic measurement noise.  The evaluation count for
\UQQAOA{} increases slightly from $45$ at $8192$ shots to $48$ at
$512$ shots; the adaptive budget rule assigns a marginally larger
batch count to noisy instances, consistent with higher effective
uncertainty.

\subsection{Analytical instantiation}
\label{sec:analytic_instantiation}

\Cref{tab:analytic_instantiation} instantiates the main bounds at the
operating scale $n{=}14$ using the isotropic proxy $\sigma_j \equiv
0.15$ and the observed 3-regular ratio $0.851$ as a proxy for $r^*$.
The table makes two points.  First, the trust-region volume ratio is
$1.44\times 10^{-4}$, so the method searches a very small fraction of
the unconstrained angle domain.  Second, the first-order landscape
bound is loose at the present scale, whereas the second-order and
noise-ordering statements are the relevant analytical guides.  The main
efficiency comparisons are also statistically significant: Wilcoxon
signed-rank tests on 48 in-distribution instances give $p<0.001$ for
\UQQAOA{} versus random initialization in evaluations and runtime, and
$p{=}0.003$ for \UQQAOA{} versus the \GNN{} point predictor in
evaluations.

\begin{table}[t]
\caption{Analytical instantiation at $n{=}14$ for 3-regular graphs
($m{=}21$, $p{=}2$). ``Bound'' denotes the theorem plug-in value under
the stated assumptions.  The learned trust region occupies a very
small fraction of the unconstrained search domain, and the
second-order and ordering guarantees are the most informative guides
at the observed operating scale.
\label{tab:analytic_instantiation}}
\begin{ruledtabular}
\scriptsize
\begin{tabular}{p{0.58\columnwidth}p{0.28\columnwidth}}
    Quantity & Bound \\
    \hline
    Lipschitz $L_\calG$ (Prop.~\ref{prop:lipschitz})
        & $84.0$ \\
    Expected-quality lower bound, Thm.~\ref{thm:landscape_lb}
        & $\ge 0.851\,C_{\max} - 38.81$ \\
    Best-of-$K$ gap (Thm.~\ref{thm:bestofk}, $K{=}3$)
        & $\le 14.55$ \\
    Volume ratio (Lem.~\ref{lem:dim_red})
        & $1.44{\times}10^{-4}$ \\
    Expected noisy lower bound, Thm.~\ref{thm:noise_robust} ($\varepsilon{=}0.01$)
        & $\ge 0.817\,C_{\max} - 36.87$ \\
    Generalization gap, Thm.~\ref{thm:generalization}
        & $\le 7.07$ \\
\end{tabular}
\end{ruledtabular}
\end{table}

% ============================================================
\section{Implications and Limitations}
\label{sec:discussion}
% ============================================================



\subsection{Origin of evaluation savings}

The method section assigns five roles: graph-conditioned prediction,
distributional alignment, representation ordering, geometric
confinement, and budget allocation.  The ablation indicates that the
largest isolated gain comes from geometric confinement.  The remaining
gain is distributed across the terms that make the covariance usable.
The Wasserstein term stabilizes distribution geometry across nearby
instances.  The contrastive term gives the representation a landscape
order before the Gaussian heads act.  The adaptive rule turns the same
covariance into a nonuniform allocation of circuit evaluations.  The
evaluation savings therefore arise from the full distributional
mechanism rather than from the center alone.

\subsection{Relation to classical trust-region methods}

\UQQAOA{} differs from classical trust-region methods~\cite{conn2000}
in that the region is predicted before optimization rather than
updated from local second-order information during optimization.
Classical schemes estimate geometry from local function data along the
trajectory.  The present scheme estimates geometry from the graph
before the first circuit query.  That difference changes the role of
learning in the algorithm.  Learning does not replace refinement.  It
supplies a center, a metric, and a budget prior to refinement.

\subsection{Interpretation as query allocation}

In low-depth \QAOA{} the dominant cost is circuit evaluations.
Restricting the search region and adapting the budget reduce this
count.  At $p{=}2$, \UQQAOA{} leads on every query-efficiency and
circuit-resource metric; at $T_{\mathrm{base}}{=}60$ the ratio gap
narrows to $0.007$.  The structural point is that the variance
head serves two purposes at once: it defines the admissible region and
it sets the amount of refinement.  Small predicted variance produces a
tight search and a short optimization.  Large predicted variance
produces a broader search and a larger budget.  The algorithm
therefore allocates quantum effort at the level of instances rather
than through a fixed global restart count.

The method is not the best absolute-quality optimizer under
unrestricted evaluation budgets.  Its contribution is near-baseline
quality at substantially lower query cost.

\subsection{Scope and limitations}
\label{sec:limitations}

The experiments are restricted to statevector simulation, $n \le 16$,
and $p{=}2$.  Weighted instances, deeper circuits, and
hardware-specific noise are not addressed.

The evidence therefore concerns a small-graph, low-depth setting.
Broader claims require new experiments.

The comparison tests whether distributional prediction improves over
random restarts, concentration-based schedules, and learned point
prediction in the same derivative-free regime; it does not cover
gradient-based methods, surrogate-model approaches, or
hardware-tailored adaptive protocols.

The diagonal covariance is axis aligned; full-covariance, low-rank,
or torus-aware extensions remain open.  Extension to higher depth,
larger systems, or hardware also remains open.

These limitations match the mechanism proved and tested in the paper.
The theory is local, the experiments are low depth, and the budget rule
is calibrated on small simulated graphs.  The present evidence is
therefore strongest in the regime in which the predictor supplies a
reliable covariance and the trust region remains computationally cheap
to use.

% ============================================================
\section{Conclusion}
\label{sec:conclusion}
% ============================================================

A graph-conditioned Gaussian predictor can organize low-depth \QAOA{}
search by supplying an initialization, a Mahalanobis trust region,
and an instance-dependent evaluation budget.  On MaxCut for four
graph families with $n{=}8$--$16$, this reduces the mean evaluation
count to $45\pm 7$, an $86.9\%$ reduction relative to random
restarts, while keeping sampled approximation ratios within $3$
percentage points of concentration-based heuristics.  At the
higher-budget operating point $T_{\mathrm{base}}{=}60$, the ratio gap
to the best absolute-quality baseline is $0.007$.  The reduction
persists across graph size, graph family, training seed, and
finite-shot simulation.  In the regime studied here, learned
uncertainty acts as an algorithmic prior that determines both the
admissible search geometry and the quantum-query budget before the
first circuit evaluation.

\section*{Author Contributions}

Molena Huynh conceived the study, developed the method, carried out the
analysis and experiments, and wrote the manuscript.

% ============================================================
\appendix
% ============================================================

\section{Implementation details}
\label{app:impl}

The training targets $\bm\theta_i^*$ are approximate local optima
obtained by multi-start Nelder--Mead optimization of the exact
statevector objective. They are not guaranteed to be global optima.
\label{app:local_optima}

\textbf{\GIN{} configuration.}
The predictor uses three \GIN{} layers with hidden dimension $64$ and
LayerNorm~\cite{ba2016} after each layer. The adjacency matrix is
row-normalized and includes self-loops. A graph representation of
dimension $128$ is obtained by concatenating mean and max pooling.
Two linear heads output $\bm\mu \in \bbR^4$ and
$\log \bm\sigma^2 \in [-5,2]^4$.

\textbf{Training.}
Optimization uses Adam~\cite{kingma2015} with learning rate $10^{-3}$
and weight decay $10^{-5}$. A cosine annealing schedule is applied over
$300$ epochs, with $150$ epochs in each of the two training phases.
Gradients are clipped at $\|\cdot\|_2 = 1$. Training is performed in
full-batch mode on $240$ graphs, with early stopping using a patience
of $50$ epochs. Total training time is approximately $2$ minutes on a
CPU.

\textbf{Loss weights.}
The Wasserstein and contrastive losses use weights
$\lambda_W = 0.1$ and $\lambda_C = 0.05$, respectively. The temperature
parameters are $\tau_W = 0.5$ and $\tau_C = 0.1$. The contrastive
threshold is
\[
\delta = \mathrm{median}\{\|\bm\theta_i^* - \bm\theta_j^*\|\}_{i<j}.
\]

\textbf{Calibration.}
The uncertainty normalization constants are
$U_{\mathrm{med}} = 1.235$ and $U_{\mathrm{iqr}} = 0.101$,
computed from the validation set of $80$ graphs.

\textbf{Baselines.}
Random initialization uses $4$ restarts with $60$ Nelder--Mead
iterations each. The concentration heuristic uses $2$ restarts from the
median training angles. The $k$-NN regressor uses $k{=}5$ with an
8-dimensional feature vector
$(\mathrm{density}, \bar{d}, \sigma_d, \bar{C}, \lambda_2,
\lambda_n, n, m)$. TQA uses a linear ramp~\cite{sack2021}. All methods
use Nelder--Mead refinement with tolerances
$x_{\mathrm{atol}} = 10^{-4}$ and $f_{\mathrm{atol}} = 10^{-6}$.

\textbf{Reproducibility.}
All experiments use random seed $42$. The implementation uses NumPy,
SciPy, NetworkX, PyTorch, and Matplotlib. Training requires
approximately $2$ minutes on a single Intel i7-12700H CPU core using
full-batch updates over $240$ graphs for $300$ epochs. Inference
requires approximately $1.2$ ms per graph for a forward pass of the
\GNN{} on CPU. The hyperparameters, training protocol, and numerical
settings required to reproduce the reported results and bound
instantiations are specified above.

\section{Query and inference cost}
\label{app:complexity}

\textbf{Training.}
Each training epoch processes $N$ graphs. A forward pass of the
\GIN{} requires $O(N L m d^2)$ operations, where $L$ is the number of
layers, $m$ is the average number of edges per graph, and $d$ is the
hidden dimension. The contrastive loss requires all-pairs similarities,
with cost $O(N^2 d)$. The Wasserstein regularizer requires pairwise
Gaussian distances, with cost $O(N^2 \cdot 2p)$. The total cost per
epoch is therefore
\[
O(N^2 d + N L m d^2).
\]
For $N{=}240$, $L{=}3$, and $d{=}64$, this cost is negligible relative
to the cost of circuit evaluations in the hybrid loop.

\textbf{Inference.}
A single forward pass of the \GIN{} requires $O(m L d^2)$ operations.
Sampling $K$ candidate parameter vectors and evaluating them on a
statevector simulator requires $O(K \cdot 2^n)$ operations. The
subsequent refinement with $T$ iterations of Nelder--Mead requires
$O(T \cdot 2^n)$. The total cost is therefore
\[
O\bigl((K + T)\,2^n\bigr).
\]

Empirically, \UQQAOA{} uses $13.1\%$ of the evaluation budget of
random initialization, corresponding to an $86.9\%$ reduction in
circuit evaluations (\cref{tab:efficiency}).

% ============================================================
\section{Best-of-$K$ bound proof}
\label{app:proof_bestofk}
% ============================================================

\begin{proof}
Let $D_k = \|\bm\theta_k - \bm\mu\|_2$. By Lipschitz continuity
(Proposition~\ref{prop:lipschitz}),
\[
F_\calG(\bm\theta_k)
\ge F_\calG(\bm\mu) - L_\calG D_k,
\]
and therefore
\[
\max_k F_\calG(\bm\theta_k)
\ge F_\calG(\bm\mu) - L_\calG \min_k D_k.
\]

Since $\bm\theta_k - \bm\mu \sim \Normal(\bm{0},\bm\Sigma)$,
\[
\bbE[D_k^2] = \mathrm{tr}(\bm\Sigma).
\]
Using $\min_k D_k^2 \le \frac{1}{K}\sum_{k=1}^K D_k^2$, we obtain
\[
\bbE[\min_k D_k]
\le \bbE\!\left[
\sqrt{\tfrac{1}{K}\sum_{k=1}^K D_k^2}
\right].
\]
By Jensen's inequality and concavity of the square root,
\[
\bbE[\min_k D_k]
\le \sqrt{\tfrac{1}{K}\sum_{k=1}^K \bbE[D_k^2]}
= \sqrt{\mathrm{tr}(\bm\Sigma)/K}.
\]

Combining the bounds yields
\[
\bbE\bigl[\max_k F_\calG(\bm\theta_k)\bigr]
\ge F_\calG(\bm\mu)
- L_\calG \sqrt{\mathrm{tr}(\bm\Sigma)/K},
\]
which proves Theorem~\ref{thm:bestofk}.
\end{proof}

% ============================================================
\section{Calibration-conditioned adaptive-allocation proof}
\label{app:proof_regret}
% ============================================================

\begin{proof}
Let $\calG$ be drawn from test distribution $\calD$. For a policy
$\pi$ allocating $K_\pi(\calG)$ samples and $T_\pi(\calG)$ polish
iterations, let $R_\pi(\calG)$ denote the per-instance regret in
\cref{eq:monotone_difficulty_model,eq:policy_regret_def}.
The total budget is $B_\pi(\calG) = K_\pi(\calG) + T_\pi(\calG)$.
The theorem does not derive a regret model from calibration alone; it
assumes the envelope \cref{eq:monotone_difficulty_model}. Taking
expectation gives
\begin{equation*}
    \mathrm{Regret}(\pi)
    \le c_0 + c_1\,\bbE_\calD\!\left[
    \frac{\sqrt{U(\calG)+\eta}}{\sqrt{K_\pi(\calG)}}
    \right].
\end{equation*}
Applying this with $\pi = \pi_{\mathrm{adapt}}$ and then using
\cref{eq:allocation_dominance} yields
\begin{equation*}
    \mathrm{Regret}(\pi_{\mathrm{adapt}})
    \le c_0 + \frac{c_1}{\sqrt{\bar K}}
    \bbE_\calD[\sqrt{U(\calG)+\eta}],
\end{equation*}
which is exactly \cref{eq:regret}. For the uniform comparator
$K_{\mathrm{unif}} \equiv \bar K$, the same envelope yields the same
right-hand side. The budget identity \cref{eq:budget_savings} is
immediate from the definition
$\alpha = \bbE_\calD[B_{\pi_{\mathrm{adapt}}}] /
\bbE_\calD[B_{\pi_{\mathrm{unif}}}]$.
Calibration and monotonicity alone do not determine the sign of
$\alpha-1$. Proposition~\ref{prop:allocation_support} shows only that
\cref{eq:K_adaptive,eq:T_adaptive} allocate more computation to larger
values of $U$; whether this yields $\alpha < 1$ depends on the
deployment distribution. The held-out experiments are consistent with
this mechanism but do not prove the dominance condition.
\end{proof}

% ============================================================
\section{Wasserstein-Lipschitz continuity proof}
\label{app:proof_wlip}
% ============================================================

\begin{proof}
Let $\phi$ denote the \GIN{} encoder and $h_\mu$, $h_\sigma$ the
prediction heads.  By the Lipschitz chain rule:
\begin{align*}
    \|\bm{g} - \bm{g}'\|_2
    &= \|\phi(\calG) - \phi(\calG')\|_2
    \le L_{\mathrm{enc}}\, d_{\mathrm{WL}}(\calG,\calG').
\end{align*}
By assumption, $d_{\mathrm{WL}}$ is a WL-compatible pseudometric
with respect to which the encoder is Lipschitz.
For the prediction heads one then has
\begin{align*}
    \|\bm\mu - \bm\mu'\|_2 &\le L_{\mathrm{head}}\|\bm{g} -
    \bm{g}'\|_2, \\
    \|\bm\sigma - \bm\sigma'\|_2 &\le L_{\mathrm{head}}\|\bm{g} -
    \bm{g}'\|_2.
\end{align*}
From the diagonal-Gaussian $W_2$ formula
(\cref{sec:wasserstein_bg}):
\begin{equation*}
    W_2^2 = \|\bm\mu - \bm\mu'\|_2^2 + \|\bm\sigma -
    \bm\sigma'\|_2^2
    \le 2 L_{\mathrm{head}}^2 \|\bm{g} - \bm{g}'\|_2^2,
\end{equation*}
and substituting the encoder bound yields
$W_2 \le \sqrt{2}\, L_{\mathrm{head}} L_{\mathrm{enc}}\,
d_{\mathrm{WL}}(\calG,\calG')$.
\end{proof}

% ============================================================
\section{Spectral-norm generalization bound proof}
\label{app:proof_gen}
% ============================================================

\begin{proof}
This argument applies standard spectral-norm Rademacher-complexity
results to the bounded-loss class induced by our fixed architecture;
no new learning-theoretic inequality is claimed. The proof uses
spectral-norm Rademacher complexity
bounds~\cite{bartlett2017,neyshabur2018,golowich2018}.
Since $\bm\mu \in \bbR^{2p}$ with angles bounded by
$\|\bm\theta^*\| \le \pi\sqrt{2p}$, and
$\bm\sigma \in [e^{-5/2}, e]^{2p}$, the $W_2^2$ loss is bounded:
$W_2^2(f_\theta(\calG), \delta_{\bm\theta^*}) \le M$ for
constant $M$ depending on the angle range and the clamping bounds.
With our clamping, $M \approx (2\pi)^2 \cdot 2p + e^2 \cdot 2p \approx 187$;
in practice, however, the empirical loss is bounded by
$M_{\mathrm{emp}} \approx 8.1$, the maximum observed $W_2^2$ on the
training set.
For an $L$-layer \GIN{} with per-layer spectral norms
$s_\ell = \|W_\ell\|_\sigma$ and input bound $B_x$,
the Rademacher complexity of the function class
$\calF = \{W_2^2(f_\theta(\cdot), \delta_{\bm\theta^*(\cdot)})\}$
satisfies~\cite{bartlett2017,golowich2018}:
\begin{equation*}
    \calR_N(\calF) \le \frac{2 B_x \prod_{\ell=1}^{L} s_\ell
    \sqrt{2L\log(2d)}}{\sqrt{N}},
\end{equation*}
where $d$ is the maximum hidden dimension.
This bound is independent of the total parameter count $D$; it depends
only on the spectral norms, depth $L$, and width $d$.
By the standard Rademacher generalization bound, with probability
$\ge 1 - \delta$:
\begin{equation*}
    \calR(f_\theta) \le \hat{\calR}_N(f_\theta)
    + 2\calR_N(\calF) + 3M\sqrt{\frac{\log(2/\delta)}{2N}},
\end{equation*}
where $\calR(f) = \bbE_\calD[W_2^2(f(\calG),
\delta_{\bm\theta^*})]$ is the population risk.
Substituting the bound above gives the result.
For our architecture ($L{=}3$, $d{=}64$,
$s_\ell \approx 1.5$ empirically,
$B_x \approx 2.1$, $N{=}240$):
\begin{equation*}
    \begin{aligned}
    2\calR_N
    &\approx 2 \times 2.1 \times 1.5^3
    \times \frac{\sqrt{6\log 128}}{\sqrt{240}} \\
    &= 2 \times 2.1 \times 3.375
    \times \frac{\sqrt{29.2}}{15.5} \\
    &\approx 4.94.
    \end{aligned}
\end{equation*}
Using the empirical loss bound $M_{\mathrm{emp}} \approx 8.1$:
\begin{equation*}
    \begin{aligned}
    3 \times 8.1 \times \sqrt{\log(40)/480}
    &\approx 24.3 \times 0.088 \\
    &\approx 2.13.
    \end{aligned}
\end{equation*}
Thus the total bound is $\hat{\calR}_N + 7.07 \approx \hat{\calR}_N + 7.1$.

The resulting estimate is still loose, but the spectral-norm scaling
($7.1$) is tighter than the covering-number bound
($\sqrt{D\log(N)/N} \approx 12.4$ for $D{=}37{,}000$) and has the
correct dependence on spectral norms.
\end{proof}

% ============================================================
\section{Expected ordering under affine depolarizing noise proof}
\label{app:proof_noise}
% ============================================================

\begin{proof}
Under global depolarizing noise with per-layer strength
$\varepsilon$, the state after $2p$ unitaries becomes
\begin{equation*}
    \rho^\varepsilon
    = (1-\nu)\ket{\bm\theta}\!\bra{\bm\theta}
    + \nu\,\frac{I}{2^n},
    \qquad
    \nu = 1-(1-\varepsilon)^{2p}.
\end{equation*}
~\cite{wang2021noise}. The noisy objective is
\begin{align*}
    F_\calG^\varepsilon(\bm\theta)
    &= \mathrm{tr}(C\,\rho^\varepsilon)
    = (1-\nu)\bra{\bm\theta}C\ket{\bm\theta}
    + \nu\,\mathrm{tr}(C/2^n) \\
    &= (1-\nu)\,F_\calG(\bm\theta) + \nu\,\bar{C},
\end{align*}
where $\bar{C} = \mathrm{tr}(C)/2^n$.
For the MaxCut cost operator
$C = \sum_{(i,j)\in\calE}\frac{1}{2}(I - Z_iZ_j)$,
each $Z_iZ_j$ is traceless, so
$\mathrm{tr}(C) = m \cdot 2^{n-1}$ and $\bar{C} = m/2$.
By Theorem~\ref{thm:landscape_lb}, for any random variable
$\Theta_{\mathrm{TR}}$ supported on $\calT_\alpha$,
\begin{equation*}
    \bbE[F_\calG(\Theta_{\mathrm{TR}})] \ge r^* C_{\max} - L_\calG
    \sqrt{\chi^2_{2p}(\alpha)\|\bm\Sigma\|_{\mathrm{op}}}.
\end{equation*}
Applying the affine map and taking expectation gives
\begin{align*}
    \bbE[F_\calG^\varepsilon(\Theta_{\mathrm{TR}})]
    &= (1-\nu)\,\bbE[F_\calG(\Theta_{\mathrm{TR}})] + \nu\,\bar{C} \\
    &\ge (1-\nu)\Bigl[
    r^* C_{\max}
    - L_\calG\sqrt{\chi^2_{2p}(\alpha)\|\bm\Sigma\|_{\mathrm{op}}}
    \Bigr] \\
    &\qquad + \nu\,\bar{C}.
\end{align*}
\begin{align*}
    &\bbE[F_\calG^\varepsilon(\Theta_{\mathrm{TR}})]
    - \bbE[F_\calG^\varepsilon(\Theta_{\mathrm{rand}})] \\
    &= (1-\nu)\Bigl(\bbE[F_\calG(\Theta_{\mathrm{TR}})]
    - \bbE[F_\calG(\Theta_{\mathrm{rand}})]\Bigr),
\end{align*}
since the $\nu\bar{C}$ terms cancel. Because $0 \le \nu < 1$ for
$\varepsilon < 1$, the sign of the expected quality difference is preserved.
The ratio in~\cref{eq:noise_amplify} follows from the affine
structure $F^\varepsilon = (1-\nu)F + \nu\bar{C}$:
\begin{align*}
    \frac{\bbE[F_\calG^\varepsilon(\Theta_{\mathrm{TR}})] - \bar{C}}
    {\bbE[F_\calG^\varepsilon(\Theta_{\mathrm{rand}})] - \bar{C}}
    &= \frac{(1-\nu)(\bbE[F_\calG(\Theta_{\mathrm{TR}})] - \bar{C})}
    {(1-\nu)(\bbE[F_\calG(\Theta_{\mathrm{rand}})] - \bar{C})} \\
    &= \frac{\bbE[F_\calG(\Theta_{\mathrm{TR}})] - \bar{C}}
    {\bbE[F_\calG(\Theta_{\mathrm{rand}})] - \bar{C}}.
\end{align*}
Whenever the denominator in \cref{eq:noise_amplify} is nonzero, this
relative excess objective is independent of the noise strength.

Finally, the noisy gradient satisfies
$\partial_j F^\varepsilon = (1-\nu)\partial_j F$, so
$\mathrm{Var}[\partial_j F^\varepsilon]
= (1-\nu)^2 \mathrm{Var}[\partial_j F]$.
In the barren-plateau regime where
$\mathrm{Var}[\partial_j F] = O(2^{-n})$~\cite{mcclean2018},
the noisy variance is doubly suppressed:
$\mathrm{Var}[\partial_j F^\varepsilon] = (1-\nu)^2 O(2^{-n})$.
Thus depolarizing noise further contracts gradients that are already
small.
\end{proof}

% ============================================================
\section{Conditional gradient anti-concentration proof}
\label{app:proof_anticonc}
% ============================================================

\begin{proof}
Let $H = -\nabla^2 F_\calG(\bm\mu)$ with
$\lambda_{\min}(H) \ge \kappa > 0$, and let
$\Lambda = \|H\|_{\mathrm{op}}$.
Assume third derivatives are bounded by $M_3$, i.e.,
$|\partial_j \partial_k \partial_\ell F_\calG(\bm\theta)| \le M_3$
for all $j,k,\ell$ and all $\bm\theta \in \calT_\alpha$.

\textbf{Step 1: Taylor expansion.}
Writing $\Delta = \bm\theta - \bm\mu$ and expanding $\partial_j F_\calG$
around $\bm\mu$:
\[
\partial_j F_\calG(\bm\theta)
= \underbrace{\partial_j F_\calG(\bm\mu)}_{=\,0 \text{ if }\bm\mu\text{ is critical}}
  \;-\; \bm e_j^\top H \Delta + R_j(\bm\theta),
\]
where the linear term uses $H = -\nabla^2 F_\calG(\bm\mu)$, and the
remainder satisfies $|R_j(\bm\theta)| \le \frac{1}{2} M_3 (2p)
\|\Delta\|_2^2 =: c_p M_3 \|\Delta\|_2^2$ by the integral form of the
Taylor remainder applied to each of the $2p$ cross-derivatives.

Denoting the linear part $g_j := -\bm e_j^\top H \Delta$, we have
$\partial_j F_\calG(\bm\theta) = g_j + R_j$.

\textbf{Step 2: Covariance of the displacement on the ellipsoid.}
For $\bm\theta$ drawn uniformly on the ellipsoid
$\calT_\alpha = \{\bm\theta : \Delta^\top \bm\Sigma^{-1} \Delta
\le \chi^2_{2p}(\alpha)\}$, the substitution
$\bm u = \bm\Sigma^{-1/2}\Delta$ maps $\calT_\alpha$ to the ball
$\{\bm u : \|\bm u\|_2^2 \le \chi^2_{2p}(\alpha)\}$.  For a uniform
distribution on a ball of radius $r$ in $\bbR^d$,
$\mathrm{Cov}[\bm u] = \frac{r^2}{d+2} I_d$.  Transforming back:
\[
\mathrm{Cov}[\Delta]
= \bm\Sigma^{1/2}\,\mathrm{Cov}[\bm u]\,\bm\Sigma^{1/2}
= \frac{\chi^2_{2p}(\alpha)}{2p+2}\,\bm\Sigma.
\]

\textbf{Step 3: Variance of the linear term.}
\[
\mathrm{Var}[g_j]
= \bm e_j^\top H\,\mathrm{Cov}[\Delta]\,H\,\bm e_j
= \frac{\chi^2_{2p}(\alpha)}{2p+2}\,
\bm e_j^\top H \bm\Sigma H \bm e_j.
\]
The matrix $H\bm\Sigma H$ is positive semidefinite, and its diagonal
entries satisfy
\[
(H\bm\Sigma H)_{jj}
= \sum_{k=1}^{2p} \sigma_k^2 (H_{jk})^2
\ge \sigma_{\min}^2 \sum_k (H_{jk})^2
= \sigma_{\min}^2 \|H \bm e_j\|_2^2
\ge \sigma_{\min}^2 \kappa^2,
\]
where the last inequality uses
$\|H \bm e_j\|_2 \ge \lambda_{\min}(H) \ge \kappa$.  Therefore
\begin{equation}
\mathrm{Var}[g_j]
\ge \frac{\kappa^2 \sigma_{\min}^2\, \chi^2_{2p}(\alpha)}{2p+2}.
\label{eq:anticonc_leading}
\end{equation}

\textbf{Step 4: Remainder correction.}
Since $\partial_j F = g_j + R_j$, the variance identity gives
\[
\mathrm{Var}[\partial_j F]
= \mathrm{Var}[g_j] + 2\,\mathrm{Cov}[g_j, R_j] + \mathrm{Var}[R_j].
\]
The covariance term is bounded using Cauchy--Schwarz:
$|{\mathrm{Cov}[g_j, R_j]}|
\le \sqrt{\mathrm{Var}[g_j]}\,\sqrt{\mathrm{Var}[R_j]}$.
For $\bm\theta \in \calT_\alpha$,
$\|\Delta\|_2 \le \sqrt{\chi^2_{2p}(\alpha)\, \sigma_{\max}^2}$,
so $|R_j| \le c_p M_3 \chi^2_{2p}(\alpha)\sigma_{\max}^2$, giving
$\mathrm{Var}[R_j] \le c_p^2 M_3^2 \chi^2_{2p}(\alpha)^2
\sigma_{\max}^4$.
The cross term is bounded by
$\Lambda M_3 C'_{p,\alpha} \sigma_{\max}^3$, where
$C'_{p,\alpha}$ absorbs the ellipsoid geometry constants.

Combining \cref{eq:anticonc_leading} with the remainder:
\[
\mathrm{Var}[\partial_j F_\calG(\bm\theta)]
\ge \frac{\kappa^2 \sigma_{\min}^2\,\chi^2_{2p}(\alpha)}{2p+2}
- C_{p,\alpha}\bigl(\Lambda M_3 \sigma_{\max}^3
+ M_3^2 \sigma_{\max}^4\bigr),
\]
which is the bound stated in Theorem~\ref{thm:anticoncentration}.  The
leading term dominates for sufficiently small $\sigma_{\max}$
relative to $\kappa/M_3$.
\end{proof}

% ============================================================
\section{Hyperparameter sensitivity}
\label{app:sensitivity}
% ============================================================

\Cref{tab:sensitivity} tests stability against modest hyperparameter
changes.

\begin{table}[h]
\caption{Hyperparameter sensitivity at $n{=}14$ (48~instances).
Default: $\lambda_W{=}0.1$, $\lambda_C{=}0.05$, $\tau_W{=}0.5$.  The
small variation across settings shows that the method is reasonably
stable to moderate loss-weight changes rather than requiring delicate
retuning.  \textbf{Bold}: best per metric.
\label{tab:sensitivity}}
\begin{ruledtabular}
\begin{tabular}{lccc}
    Setting & $\rho$ & Evals
    & Med.\ ms \\
    \hline
    Default & $\mathbf{0.770}$ & $\mathbf{45}$ & $\mathbf{69}$ \\
    $\lambda_W = 0.01$ & $0.742$ & $47$ & $76$ \\
    $\lambda_W = 0.5$  & $0.758$ & $46$ & $73$ \\
    $\lambda_C = 0.0$  & $0.751$ & $48$ & $77$ \\
    $\lambda_C = 0.2$  & $0.763$ & $46$ & $72$ \\
    $\tau_W = 0.2$     & $0.761$ & $46$ & $73$ \\
    $\tau_W = 1.0$     & $0.749$ & $47$ & $75$ \\
\end{tabular}
\end{ruledtabular}
\end{table}

The default setting gives the best combination.  Variation is
$\pm 0.028$ in $\rho$ and $\pm 3$ evaluations.  The narrow range
confirms that the method is not sensitive to the specific loss
weights; moderate perturbations produce results within the
statistical uncertainty of the primary experiments.

% ============================================================
\section{Statistical significance}
\label{app:significance}
% ============================================================

Wilcoxon signed-rank tests on 48 in-distribution instances confirm
the main evaluation-count and runtime comparisons: \UQQAOA{} vs.\ random in
evaluations ($p < 0.001$, effect size $r = 0.86$), \UQQAOA{} vs.\ random
in runtime ($p < 0.001$, $r = 0.84$), and \UQQAOA{} vs.\ \GNN{} point in
evaluations ($p{=}0.003$, $r = 0.42$). For calibration,
Spearman $\rho{=}0.770$ with $p < 10^{-47}$ across 240 instances.

% ============================================================
\section*{Reproducibility Statement}
\label{app:reproducibility}
% ============================================================
\begin{itemize}[leftmargin=*,nosep]
\item \textbf{Code and inputs:} The manuscript and appendix provide the
    numerical inputs used for the reported bound instantiations.
\item \textbf{Hyperparameters:} All training and optimization
      hyperparameters are listed in Appendix~\ref{app:impl}.
\item \textbf{Random seeds:} Every reported result is averaged over
      five independent training seeds; per-seed statistics are
      provided in \cref{tab:multiseed}.
\item \textbf{Data:} All experiments use synthetic graph instances
      generated with documented parameters
      (family, size, and edge-probability ranges specified in
      \cref{sec:setup}).
\end{itemize}

\section*{Data Availability}

The data that support the findings of this study are available from
the corresponding author upon reasonable request.

\hypersetup{hidelinks}
\bibliographystyle{apsrev4-2}
\bibliography{refs}

\end{document}
